\section{Solutions and fluctuations}
\label{ch:solutions}

This chapter treats the classical solutions of the cubic action and the fluctuations around them: level truncation and the equations a truncated solution violates, the Hessian of the cubic action and its congruence with the Hessian of the eliminated action, the alternative elimination of the string field, the degenerate branch of the auxiliary equation, and a mean-field reading of the auxiliary field.
The paper keeps only a short summary of these results, in the paragraphs on the Hessian, on level truncation and on the mean field of its section of further results.
The full text below is the section on classical solutions and the appendix with the proof of the truncation lemma of an earlier version of the paper, reproduced as they stood.
Only this write-up contains the analysis of the dilaton direction of the tachyon--dilaton truncation at quintic order, in \cref{sec:sol:dilaton-d5}, written from its derivation record.
The quartic vertex at level $(1, 1)$ is that of the companion paper on $\group{SL}(2, \C)$ vertices~\cite{Erbin:2026:SL2CVerticesClosed}; its numerical values, and the critical points of the tachyon--dilaton potential computed from them, are not reproduced here.

\paragraph{Status of the results.}

The results of the chapter and their status, in the order of the text:
\begin{itemize}
  \setlength{\itemsep}{0.5ex}

  \item Level truncation of the auxiliary field to $\mathcal U^{\wedge 2}$, exact elimination at quartic order and the matching through quintic and sextic order with the truncation-consistent reduced vertex:
  \status{Proven}\record{057-finite-field-truncation (R61)}

  \item Which equations of the untruncated theory a truncated solution solves, and the selection rule that protects the directions towards the non-universal states (the truncation lemma):
  \status{mixed}\record{057-finite-field-truncation, part B (R63)}

  \item The one-leg universality selection rule at zero momentum, assumed in the text below and established in the literature:
  \status{Proven}\footnote{Ledger entry \recordref{R88}; a cited result (Yang and Zwiebach, after Sen), with no derivation record of its own.}

  \item The joint Hessian of the cubic action is congruent to the Hessian of the eliminated action plus the auxiliary block (Schur identity):
  \status{Proven -- Verified (finite model)}\record{060-mean-field-fluctuations (R65)}

  \item The Hessian of the cubic action is affine in the background; fluctuation counts and zero modes from the low vertices only:
  \status{Proven; C2.3 Verified (finite model)}\record{066-hessian-affine-background (R71)}

  \item Level-$(1, 1)$ instance: exact elimination at quartic order, the indefinite auxiliary form and the role of the dilaton:
  \status{Verified}\record{059-tachyon-dilaton-quartic (R64)}

  \item Level-$(1, 1)$ instance: the cubic action at zero momentum, its indefinite auxiliary block and the contour of the Gaussian integral:
  \status{Verified}\record{064-hs-action-level11 (R68)}

  \item Level-$(1, 1)$ instance: the congruence of the joint Hessian at the critical points, and the wrong count of the frozen auxiliary field:
  \status{Verified}\record{063-mean-field-level11 (R72)}

  \item Integrating out the string field at tree level: the operator of the auxiliary field, the branch structure of the $\Psi$-equation, the tree-level functional of $\Sigma_2$:
  \status{Conditional(PH1, PH5) -- Proven}\record{065-sigma-effective-action (R69)}

  \item Exclusion, in a finite model at quintic order with a kernel of dimension one: the degenerate branch carries no generic joint critical point:
  \status{Proven -- Verified (finite model)}\record{060-mean-field-fluctuations, part 3 (R67)}

  \item Mean field: the one-loop potential of the auxiliary field, the gap equation, and the obstacles to a condensate of BCS type:
  \status{Proven -- Verified (finite model); Argued (obstacles)}\record{069-mean-field-decomposition (R74)}

  \item Mean field at level $(1, 1)$: no minimum of the gap equation on the real slice; the convergent saddles are the displaced classical points:
  \status{Proven (reduction, limit points, identity); Verified (numerical) (root set, inertia, energies, contour patterns)}\record{070-mean-field-level11-gap (R77)}

  \item Exclusion, in the five-field truncation: no flat dilaton direction at order $d^5$ (\cref{sec:sol:dilaton-d5}):
  \status{Conditional(PH27, SL: PH5) -- mixed}\record{095-d5-five-field-truncation (R89)}
\end{itemize}

\paragraph{The section on classical solutions.}

The next block is the section on classical solutions and fluctuations of the earlier version of the paper, in the paper's notation ($\vertex{n}$ for the vertices, $\sprod{2}{2}$ for the quartic form read as an operator, $\caction$ for the cubic action, $V = - S$ on the potential side); its headings are demoted by one level.
Two remarks on its status.
First, the text uses the universality rule~\eqref{p:eq:sol:universality}.
The rule is established at genus zero and zero momentum, on the Fock space of zero momentum: a vertex given as a moduli integral of a sphere correlator of conformally transformed states with $b$-ghost insertions vanishes when exactly one leg has its matter part outside the matter identity module and all the others inside it, whatever the ghost parts~\cite{Yang:2005:ClosedStringTachyon,Sen:1999:UniversalityTachyonPotential}.
It holds for the cubic vertex and for the quartic vertex of the $\group{SL}(2, \C)$ family in any admissible section; at five points and more it rests on the family-independent Ward argument, since no quintic vertex of that family exists yet.
The kinetic term obeys the same rule at every level and ghost number, and at nonzero momentum the same couplings are of order $k^2$.
The rule does not license dropping the exchanges outside the module in limits $k \to 0$ taken through massless poles; it is used here at fixed zero momentum only.
Second, in the level-$(1, 1)$ instance the statements that the couplings lie in the window of indefiniteness and that the frozen auxiliary field gives the wrong count are read off the companion paper's values, which we do not reproduce.

\subsection{Classical solutions and fluctuations}
\label{p:sec:solutions}

At quartic order the cubic action and the original one have the same classical solutions, since eliminating the auxiliary field is exact and gives back the original action.
At higher orders the solutions of the cubic action at which the auxiliary field can be eliminated, that is, where its quadratic operator is invertible, are the critical points of the eliminated action, which agrees with the original truncated action through the order kept; the configurations where it cannot are treated with the degenerate branch below.
The difference lies in the fluctuation problem around a solution: the second variation of a cubic action is affine in the background, while that of the original action is not.
We make the two statements precise, first in the level truncations in which solutions are actually computed, then for the Hessian, and finally for the two alternative eliminations, of $\Sigma_2$ and of $\Psi$.
The statements are algebraic: they use the polynomial form of the action and linear algebra.
They hold for the string vertices wherever the inverses they involve exist, under the assumptions of \cref{p:sec:hs}: the quartic form invertible on the quotient, the quintic and sextic vertices vanishing on its kernel.

\subsubsection{Level truncations}
\label{p:sec:sol:level-truncations}


Solutions of closed string field theory are found numerically by level truncation, restricting the string field to a finite-dimensional subspace, typically the states up to a given level in Siegel gauge~\cite{Yang:2005:ClosedStringTachyon}.
Let $\mathcal U \subset \Hs$ be a finite-dimensional subspace, and let $\Hs = \mathcal U \oplus \mathcal U^c$ be a decomposition in which the kinetic term does not mix the two summands:
\begin{equation}
  \label{p:eq:sol:kinetic-block}
  \pair{u}{Q_B x}
    = 0,
  \qquad
  u \in \mathcal U,
  \quad
  x \in \mathcal U^c,
\end{equation}
which is the case when $\mathcal U$ is a union of levels and $\mathcal U^c$ the remaining levels, since $Q_B$ commutes with $L_0^+$ and the pairing couples only states of equal level.
We extend the level truncation of the string field to the auxiliary field by restricting it to $\mathcal U^{\wedge 2}$, the symmetric square of $\mathcal U$:
\begin{equation}
  \label{p:eq:sol:truncation}
  \Psi
    \in \mathcal U,
  \qquad
  \Sigma_2
    \in \mathcal U^{\wedge 2},
  \qquad
  \caction[\le N]^{[\mathcal U]}
    := \caction[\le N] \big|_{\mathcal U \times \mathcal U^{\wedge 2}},
\end{equation}
with the vertices evaluated as in the untruncated theory, and we assume that the quartic form has no kernel on $\mathcal U^{\wedge 2}$, so that no quotient is needed there.
The choice $\mathcal U^{\wedge 2}$ is the smallest one that works: it contains $\Psi^2$ for every $\Psi \in \mathcal U$, so the solution $\Sigma_2^{(2)*} = \Psi^2$ of \cref{p:sec:hs:quartic} is representable in the truncation.
An auxiliary space which does not contain $\mathcal U^{\wedge 2}$ fails already at quartic order: the square can no longer be completed, and the eliminated action differs from the quartic truncation for generic couplings.

Eliminating $\Sigma_2$ from~\eqref{p:eq:sol:truncation} gives an action for $\Psi \in \mathcal U$ which we compare with the level truncation of the original theory, $S_{\le N} \big|_{\mathcal U}$.
At quartic order the two agree exactly, since the solution $\Sigma_2 = \Psi^2$ lies in $\mathcal U^{\wedge 2}$ and the square completes inside the truncation.
At quintic order the operator of $\Sigma_2$ is the dressed operator $\kin{2}{\Psi}{3}$ of \cref{p:sec:hs:quintic} restricted to $\mathcal U^{\wedge 2}$, and its inverse is a rational function of $\Psi$.
The eliminated action then agrees with $S_{\le 5} \big|_{\mathcal U}$ through $O(\gs^3)$, with the same exchange term at $O(\gs^4)$ as in the untruncated theory, up to the inverse used.
Writing $\Sigma_2^{(3)*[\mathcal U]}$ for the solution of the truncated auxiliary equation on $\mathcal U^{\wedge 2}$, we get:
\begin{equation}
  \label{p:eq:sol:matching}
  \begin{gathered}
    \caction[\le 4]^{[\mathcal U]} \Big|_{\Sigma_2 = \Psi^2}
      = S_{\le 4} \big|_{\mathcal U},
    \\
    \caction[\le 5]^{[\mathcal U]} \Big|_{\Sigma_2 = \Sigma_2^{(3)*[\mathcal U]}}
      = S_{\le 5} \big|_{\mathcal U}
        + \frac{\gs^4}{600} \, \vertex{5 \cdot 5}^{[\mathcal U]}(\Psi^6)
        + O(\gs^5).
  \end{gathered}
\end{equation}
Here $\vertex{5 \cdot 5}^{[\mathcal U]}$ is the composite vertex of \cref{p:sec:hs:quintic} built with $(\sprod{2}{2}^{[\mathcal U]})^{-1}$, the inverse of the quartic form on $\mathcal U^{\wedge 2}$, in place of $\sprod{2}{2}^{-1}$.
At sextic order the reduced vertex $\dvertex{6} = \vertex{6} - \tfrac{6}{5} \vertex{5 \cdot 5}$ of \cref{p:sec:hs:sextic} is built with the inverse of $\sprod{2}{2}$ on the whole of $\Hn{2}$, which is not computable inside a truncation.
If it is nevertheless used, with that inverse computed outside the truncation, the eliminated action agrees with $S_{\le 6} \big|_{\mathcal U}$ through $O(\gs^3)$.
The $O(\gs^4)$ difference is quadratic in the part of the source $\sprod{3}{2}(\Psi^3)$ that the truncated auxiliary space cannot absorb, with coefficient the inverse of the Schur complement (\cref{p:sec:hs:sextic:kernel}) of the block of $\vertex{4}$ on $\mathcal U^{\wedge 2}$~\eqref{p:eq:trunc:schur}.
That part vanishes exactly when $\sprod{2}{2}^{-1} \circ \sprod{3}{2}(\Psi^3)$ lies in $\mathcal U^{\wedge 2}$.
The vertex to use in a truncation is instead the truncation-consistent reduced vertex
\begin{equation}
  \label{p:eq:sol:W6-truncated}
  \dvertex{6}^{[\mathcal U]}
    := \vertex{6} - \frac{6}{5} \, \vertex{5 \cdot 5}^{[\mathcal U]},
\end{equation}
built from the same data as the rest of the truncated action.
With it the eliminated action equals $S_{\le 6} \big|_{\mathcal U} + O(\gs^5)$, without enlarging the auxiliary space: the truncated exchange term of the quintic order is cancelled by the truncated composite vertex inside $\dvertex{6}^{[\mathcal U]}$, as in the untruncated theory.
The price is that this truncated action is not the restriction of the full one.
Its solutions violate the equation of $\Sigma_2$ of the full cubic action along $\mathcal U^{\wedge 2}$ itself at relative order $\gs^2$.
This violation and the one along the pairs with a leg outside $\mathcal U$ cancel in the value of the action through order $\gs^4$: this is the matching with $S_{\le 6} \big|_{\mathcal U}$ just stated, proved in \cref{p:app:truncation}.
These matchings are statements about Taylor coefficients in $\gs$ at fixed $\Psi$, that is, about the terms through $\Psi^5$ or $\Psi^6$ of the action.
At a tachyon vacuum $\gs \Psi = O(1)$, every term is of the same size, and the matchings carry no accuracy statement there.

For a numerical solution the relevant question is different: a solution of the truncated equations of motion is not a solution of the untruncated ones, and we want to know which equations it violates and by how much.
For the original theory the answer is the standard level-truncation error, the $\Psi$-equation along $\mathcal U^c$.
For the cubic theory there is a second kind of error, along the auxiliary directions with a leg in $\mathcal U^c$.
A selection rule, Sen's universality~\cite{Sen:1999:UniversalityTachyonPotential}, removes both along the directions that point towards the non-universal states, at the orders stated below (\cref{p:app:truncation}).
At zero momentum, a vertex with all legs universal (matter part in the identity Virasoro module) except exactly one vanishes, by a Ward-identity argument: it is an integral over moduli of sphere correlators with a single matter leg outside the identity module.
The same argument underlies the level-truncation studies of the closed-string tachyon vacuum~\cite{Yang:2005:ClosedStringTachyon}.
The rule is established at genus zero and zero momentum, on the Fock space of zero momentum, for any vertex given as a moduli integral of a sphere correlator of conformally transformed states with $b$-ghost insertions, and it holds for the cubic and quartic vertices of the $\group{SL}(2, \C)$ family in any admissible section; at five points and more it rests on this family-independent Ward argument, since no quintic vertex of that family exists yet.\footnote{Ledger entry \recordref{R88}; a cited result (Yang and Zwiebach, after Sen), with no derivation record of its own.}
We use it in the following form, for the vertices this section uses:
\begin{equation}
  \label{p:eq:sol:universality}
  \vertex{n}(u_1, \ldots, u_{n-1}, x)
    = 0,
  \qquad
  u_i \in \mathcal U,
  \quad
  x \in \mathcal N,
  \quad
  3 \le n \le 6,
\end{equation}
for $\mathcal U$ a subspace of zero-momentum universal states in Siegel gauge, such as the tachyon and the ghost dilaton, and $\mathcal N$ the non-universal states; nonzero-momentum states satisfy~\eqref{p:eq:sol:universality} by momentum conservation.
The kinetic term does not mix $\mathcal U$ and $\mathcal N$ either, since $Q_B$ preserves the matter Virasoro module and the pairing does not mix modules.
From here on $\mathcal U$ is such a space of universal states.
  Then~\eqref{p:eq:sol:kinetic-block} holds towards $\mathcal N$ by this argument, and towards the universal states of other levels by the level argument below~\eqref{p:eq:sol:kinetic-block}.
  The universal states of the same levels are discussed at the end of this subsection.

Under these assumptions, let $(\Psi^*, \Sigma_2^*)$ solve the equations of motion of the truncated cubic action~\eqref{p:eq:sol:truncation} on the perturbative branch, the one connected to $\Sigma_2 = \Psi^2$.
The $\Psi$-equation along $\mathcal N$ then vanishes identically in the fields, at every order considered here.
At quartic order this is visible term by term: for every $x \in \mathcal N$,
\begin{equation}
  \label{p:eq:sol:lemma}
  \pair{x}{\mathcal E^{(2)}_\Psi(\Psi^*, \Sigma_2^*)}
    = \pair{x}{Q_B \Psi^*}
      + \frac{\gs}{2} \, \vertex{3}(x, \Psi^{* 2})
      + \frac{\gs^2}{6} \, \vertex{4}(x, \Sigma_2^*, \Psi^*)
    = 0,
\end{equation}
the first term by~\eqref{p:eq:sol:kinetic-block} and the other two by~\eqref{p:eq:sol:universality}, since $\Sigma_2^* \in \mathcal U^{\wedge 2}$ has both strands in $\mathcal U$.
Since the $\Sigma_2$-equation holds exactly at quartic order, a solution of the quartic truncation solves every equation of the quartic-order cubic theory except the $\Psi$-equation along the universal states outside $\mathcal U$.

At quintic order the $\Sigma_2$-equation holds identically along the pairs of a state of $\mathcal U$ and a non-universal state.
It is violated along the pairs of two non-universal states and, in general, along the pairs with a universal state outside $\mathcal U$, at relative order $\gs$, sourced by the part of $\sprod{3}{2}(\Psi^{* 3})$ that $\mathcal U^{\wedge 2}$ cannot absorb.
At sextic order it is in general violated along the pairs of one state of $\mathcal U$ and one non-universal state as well, at relative order $\gs^2$, through the composite vertex.
With~\eqref{p:eq:sol:W6-truncated}, so is $\mathcal U^{\wedge 2}$ itself, at the same relative order.
Solving these auxiliary equations at fixed $\Psi^*$ changes the action by the term of order $\gs^4 \Psi^6$ found above, which~\eqref{p:eq:sol:W6-truncated} cancels.
Along the universal states outside $\mathcal U$ the rule says nothing: for a finite $\mathcal U$ such as the tachyon and the ghost dilaton below, the error along the universal states of higher level is the ordinary level-truncation error.
Same-level universal states outside $\mathcal U$ – outside Siegel gauge, or in Siegel gauge but not kept in $\mathcal U$ – are not removed by the rule, the non-universal ones being in $\mathcal N$, and~\eqref{p:eq:sol:kinetic-block} is not established for them.
Their equations are the out-of-gauge equations of level truncation, which we do not examine.

\subsubsection{The Hessian of the cubic action}
\label{p:sec:sol:hessian}


Consider a background $(\Psi^*, \Sigma_2^*)$ and the fluctuations around it:
\begin{equation}
  \label{p:eq:sol:fluctuations}
  \Psi
    = \Psi^* + \psi,
  \qquad
  \Sigma_2
    = \Sigma_2^* + \sigma_2.
\end{equation}
The Hessian $\Hess \caction$ is the second variation of the action at the background, a symmetric bilinear form on the fluctuations with three blocks.
Since the action is cubic, every block is affine in the background.
The second derivative of a cubic monomial is linear in the fields and that of a quadratic one is constant, so each background-dependent term of a block is one vertex contracted once with one background field.
Through quintic order the blocks are:
\begin{subequations}
  \label{p:eq:sol:hessian-blocks}
  \begin{align}
    \label{p:eq:sol:H-psi-psi}
    H_{\Psi \Psi}(\psi, \psi')
      &
      = \pair{\psi}{Q_B \psi'}
        + \gs \, \vertex{3}(\Psi^*, \psi, \psi')
        + \frac{\gs^2}{6} \, \vertex{4}(\Sigma_2^*, \psi, \psi'),
    \\
    \label{p:eq:sol:H-sigma-psi}
    H_{\Sigma_2 \Psi}(\sigma_2, \psi)
      &
      = \frac{\gs^2}{6} \, \vertex{4}(\sigma_2, \Psi^*, \psi)
        + \frac{\gs^3}{60} \, \vertex{5}(\sigma_2, \Sigma_2^*, \psi),
    \\
    \label{p:eq:sol:H-sigma-sigma}
    H_{\Sigma_2 \Sigma_2}(\sigma_2, \sigma_2')
      &
      = - \frac{\gs^2}{12} \, \Big[
          \vertex{4}(\sigma_2, \sigma_2')
          - \frac{\gs}{5} \, \vertex{5}(\sigma_2, \sigma_2', \Psi^*)
        \Big]
      = - \frac{\gs^2}{12} \, \pair[2]{\sigma_2}{\kin{2}{\Psi^*}{3}(\sigma_2')},
  \end{align}
\end{subequations}
where the auxiliary block is the dressed kinetic operator of \cref{p:sec:hs:quintic} evaluated in the background.

The $\Psi \Psi$ block receives no further term at any order, since the only cubic monomials containing $\Psi$ twice are $\Psi^3$ and $\Sigma_2 \Psi^2$.
The $\Sigma_2 \Psi$ block is complete at quintic order.
  When the sextic term is written with $\Sigma_2$ only (\cref{p:sec:hs:sextic}), the $\Sigma_2 \Sigma_2$ block is complete at sextic order, where the reduced sextic vertex contracted with $\Sigma_2^*$ enters; with $\Sigma_3$ the sextic terms go instead to its blocks and to the $\Sigma_2 \Psi$ block.
  Each higher auxiliary field $\Sigma_k$, $k \ge 3$, adds blocks of the same kind.
Its diagonal block is constant at the order at which the field enters and linear in $\Psi^*$ one order later, and it couples to $\Sigma_{k \pm 1}$ and to $\Psi$ only.

We write $\phi := (\psi, \sigma_2)$ for the joint fluctuation.
The fluctuation operator $\mathcal T$ is obtained from the Hessian by the kinetic pairing $\pair{\cdot}{\cdot} \oplus \pair[2]{\cdot}{\cdot}$:
\begin{equation}
  \label{p:eq:sol:fluctuation-operator}
  \Hess \caction(\phi, \phi')
    =: \pair{\psi}{(\mathcal T \phi')_\Psi}
      + \pair[2]{\sigma_2}{(\mathcal T \phi')_{\Sigma_2}},
\end{equation}
so that $\mathcal T$ is one affine matrix built from the vertices through $\vertex{5}$ and the background.

The joint Hessian is not the fluctuation operator of the eliminated theory, but the two are related by a change of variables.
Wherever the auxiliary block is invertible, that is, wherever $\kin{2}{\Psi^*}{3}$ is invertible on the quotient, we complete the square in the auxiliary fluctuation.
In the shifted auxiliary fluctuation
\begin{equation}
  \label{p:eq:sol:shift}
  x_2
    := \sigma_2 + H_{\Sigma_2 \Sigma_2}^{-1} H_{\Sigma_2 \Psi} \, \psi,
\end{equation}
the Hessian becomes block diagonal:
\begin{equation}
  \label{p:eq:sol:congruence}
  \begin{gathered}
    \Hess \caction
      = L^{\mathsf T} \,
        \big( H_{\rm eff} \oplus H_{\Sigma_2 \Sigma_2} \big) \,
        L,
    \qquad
    L
      = \begin{pmatrix}
          1 & 0
          \\
          H_{\Sigma_2 \Sigma_2}^{-1} H_{\Sigma_2 \Psi} & 1
        \end{pmatrix},
    \\
    H_{\rm eff}
      := H_{\Psi \Psi}
        - H_{\Sigma_2 \Psi}^{\mathsf T} \, H_{\Sigma_2 \Sigma_2}^{-1} \, H_{\Sigma_2 \Psi},
  \end{gathered}
\end{equation}
where $H_{\rm eff}$ is the Schur complement of the auxiliary block and $L$ is unipotent.

Two forms related by an invertible $L$ have the same inertia, the numbers of positive, negative and zero directions (Sylvester's law), and the same kernel up to $L$.
We write $\mathrm{In}$ for this triple.
They do not, in general, have the same eigenvalues with respect to the kinetic pairing, since $L$ does not preserve that pairing.
Hence, in any finite truncation, the counts add:
\begin{equation}
  \label{p:eq:sol:inertia}
  \begin{gathered}
    \mathrm{In}(\Hess \caction)
      = \mathrm{In}(H_{\rm eff}) + \mathrm{In}(H_{\Sigma_2 \Sigma_2}),
    \\
    \ker \Hess \caction
      = \big\{ (\psi, \, - H_{\Sigma_2 \Sigma_2}^{-1} H_{\Sigma_2 \Psi} \, \psi) : \psi \in \ker H_{\rm eff} \big\},
  \end{gathered}
\end{equation}
and the determinant factorizes, $\det \Hess \caction = \det H_{\Sigma_2 \Sigma_2} \, \det H_{\rm eff}$.

On the auxiliary shell, $\Sigma_2^* = \Sigma_2^{(3)*}(\Psi^*)$, the shift $x_2$ is the linearization of $\Sigma_2 - \Sigma_2^{(3)*}(\Psi)$, and $H_{\rm eff}$ is exactly the Hessian of the eliminated action.
There, $- H_{\Sigma_2 \Sigma_2}^{-1} H_{\Sigma_2 \Psi} \, \psi$ is the derivative of $\Sigma_2^{(3)*}(\Psi)$ along $\psi$.
This is the envelope theorem at second order: differentiating the eliminated action $\caction\big(\Psi, \Sigma_2^{(3)*}(\Psi)\big)$ twice and using the auxiliary equation reproduces the Schur complement.
The identity holds on and off the $\Psi$-shell, at every order, on any branch of the auxiliary equation where the auxiliary block is invertible; at sextic order this excludes the fold discussed below.
Through the matchings of \cref{p:sec:hs}, $H_{\rm eff}$ equals the Hessian of the original truncated action exactly at quartic order and through $O(\gs^3)$ at quintic order, the $O(\gs^4)$ remainder being the Hessian of the exchange term.
At sextic order the two agree through $O(\gs^4)$.
The zero mode of the cubic theory is the zero mode of the eliminated theory dressed by the fluctuation of the auxiliary field which follows the solution.
The naive lift $(\psi, 0)$ is a zero mode only if $H_{\Sigma_2 \Psi} \psi = 0$.

At quartic order the split can be read on the vertex itself.
With $\Sigma_2^* = \Psi^{* 2}$ the Schur term is $\tfrac{\gs^2}{3} \vertex{4}(\Psi^{* 2}, \psi, \psi)$, and we get:
\begin{equation}
  \label{p:eq:sol:channels}
  \begin{aligned}
    H_{\rm eff}(\psi, \psi)
      &
      = \pair{\psi}{Q_B \psi}
        + \gs \, \vertex{3}(\Psi^*, \psi, \psi)
        + \underbrace{\frac{\gs^2}{6} \, \vertex{4}(\Psi^{* 2}, \psi, \psi)}_{\text{direct channel}}
        + \underbrace{\frac{\gs^2}{3} \, \vertex{4}(\Psi^{* 2}, \psi, \psi)}_{\sigma_2 \text{ exchange}}
      \\
      &
      = \pair{\psi}{Q_{\Psi^*} \psi},
  \end{aligned}
\end{equation}
with $Q_{\Psi^*}$ the shifted BRST operator of the original theory around $\Psi^*$ (\cref{p:sec:gauge}), here at quartic order~\cite{Zwiebach:1993:ClosedStringField}.
This is the split into the direct channel and the exchange found for the BRST operator around the background in~\eqref{p:eq:gauge:direct-channel}, here read on the Hessian.
At quintic order the $\Psi \Psi$ block carries one fifth of the quintic term, through the $O(\gs)$ correction of $\Sigma_2^*$, and the exchange of $\sigma_2$ the remaining four fifths.

This is also where the two descriptions differ in kind.
The original action at quartic order is quartic in $\Psi$, so its Hessian is quadratic in the background.
The Hessian of the cubic action is affine, and the quadratic dependence is produced by the elimination, through the mixed block $H_{\Sigma_2 \Psi}$, linear in $\Psi^*$, which enters the Schur complement twice.

At higher orders the elimination also produces singularities that the cubic theory does not have.
At sextic order, when the sextic term is written with $\Sigma_2$ only, without $\Sigma_3$ (\cref{p:sec:hs:sextic}), the auxiliary equation is quadratic in $\Sigma_2$ and its solution has two sheets which meet where $H_{\Sigma_2 \Sigma_2}$ degenerates, a fold.
There the Hessian of the eliminated action diverges as the inverse square root of the distance of $\Psi^*$ to the fold, approached along either sheet, while $\Hess \caction$, being affine, is finite and continuous, with the Schur complement simply undefined.
In the one-field model of \cref{p:app:finite-models}, the joint Hessian at the fold has a determinant that does not vanish for generic couplings, with one positive and one negative direction.
In that model the two summands of~\eqref{p:eq:sol:inertia} therefore exchange one unit between the sheets meeting at the fold, and the total count stays regular.
The fold is a property of the elimination, not of the cubic theory.

Two cautions are in order.
The kinetic pairing is indefinite, so that the fluctuation operator $\mathcal T$, although self-adjoint with respect to it, is in general not self-adjoint for a positive inner product: its eigenvalues need not be real, and their signs are not a count of unstable directions.
The invariant quantities, which~\eqref{p:eq:sol:congruence} controls, are the inertia and the kernel.
Once momentum is restored, a mass shell is a momentum at which the fluctuation operator has a zero mode.
Then~\eqref{p:eq:sol:inertia} says that at momenta where the auxiliary block is invertible the mass shells of the cubic theory around the background are those of $H_{\rm eff}$.
We know of nothing that excludes a zero of $\det H_{\Sigma_2 \Sigma_2}(k)$ at some $k \neq 0$: the contact vertices give the block no pole, but its determinant can vanish at isolated momenta.
Such a zero would be a mass shell of the cubic theory carried by the auxiliary field and absent from $H_{\rm eff}$.

\subparagraph{The level-\texorpdfstring{$(1, 1)$}{(1,1)} instance.}


We illustrate the structure on the truncation to the tachyon $T$ and the ghost dilaton $D$ at level $(1, 1)$ and zero momentum, with the quartic vertex of the companion paper~\cite{Erbin:2026:SL2CVerticesClosed} (the auxiliary field restricted to $\{T, D\}^{\wedge 2}$ is not a consistent truncation, since $\vertex{4}$ couples $T \otimes T$ and $D \otimes D$ to pairs of matter states).
Levels are counted from the tachyon in the left and right sectors, so that level $(1, 1)$ is $L_0 = \bar L_0 = 0$, the level of the dilaton, and the truncation keeps the states up to it.
Here $D = (c_1 c_{-1} - \bar c_1 \bar c_{-1}) \ket{0}$, with no matter oscillator, is the same vector as in the companion paper, and $\braket{D | c_0^- c_0^+ | D} = 1$ in our conventions.
The string field has two components and the auxiliary field three:
\begin{equation}
  \label{p:eq:sol:level11}
  \begin{gathered}
  \Psi
    = t \, T + d \, D,
  \\
  \Sigma_2
    = \sigma_{TT} \, T \otimes T
      + \sigma_{TD} \, \tfrac{1}{2} (T \otimes D + D \otimes T)
      + \sigma_{DD} \, D \otimes D
    \in \mathcal U^{\wedge 2},
  \quad
  \mathcal U
    = \Span(T, D),
  \end{gathered}
\end{equation}
so that the cubic action is a function of five real variables, cubic in $(t, d)$ and quadratic in $\sigma$.
At quartic order the elimination is exact, $\Sigma_2^* = \Psi^2$, that is, $(\sigma_{TT}, \sigma_{TD}, \sigma_{DD}) = (t^2, 2 t d, d^2)$.
The eliminated action is the ordinary $2$-field quartic potential of the level-truncated original theory, for this vertex family: nothing specific to the cubic formulation enters before the quintic vertex.
Its critical points, their values and the comparison with the literature are the subject of the companion paper, and we only use their structure here.

At this order the cubic formulation adds the auxiliary block.
The quartic form $\vertex{4}$ restricted to $\mathcal U^{\wedge 2}$ is indefinite.
On the potential side, $V := - S \big|_{k = 0}$ as in \cref{p:app:conventions}, the block $\tfrac{\gs^2}{12} \vertex{4}$ has exactly one negative direction, and the tachyon--dilaton pair $T \otimes D$ is itself a negative direction, since $\vertex{4}(T^2, D^2) < 0$.
This holds whenever $\vertex{4}(T^4)$ and $\vertex{4}(D^4)$ are positive and $\vertex{4}(T^2, D^2)$ lies between $- \sqrt{\vertex{4}(T^4) \, \vertex{4}(D^4)}$ and zero, whatever the two couplings with an odd number of dilatons.
The Gaussian integral over $\Sigma_2$, formally exact whatever the signature (\cref{p:sec:hs:quartic}), is then defined by rotating this one direction to the imaginary axis.
The result is the exponential of the quartic term times a field-independent constant, whose phase alone depends on the orientation of the rotated axis, and the propagator of $\Sigma_2$ is unchanged.
The negative direction is the auxiliary analogue of the conformal-factor rotation~\cite{Gibbons:1978:PathIntegralsIndefiniteness}, not an instability: at zero momentum the block is a constant, algebraic direction of the quadratic form, with no $L_0^+$ and no derivative in its equation.

Around the perturbative vacuum and around each regular critical point of the $2$-field potential, the joint Hessian is the $5 \times 5$ matrix~\eqref{p:eq:sol:hessian-blocks} with $\vertex{5} = 0$.
The congruence~\eqref{p:eq:sol:congruence} holds there with $H_{\rm eff}$ equal to minus the Hessian of the $2$-field potential $V$.
The auxiliary sector adds its constant inertia at every background, one negative direction on the potential side.
It is orthogonal, with respect to the Hessian, to the tangent of the auxiliary shell $(\psi, \partial_\Psi \Sigma_2^{(2)*}(\Psi^*) \, \psi)$, on which the Hessian is that of the eliminated theory.
Everything the zero-momentum truncation says about a critical point therefore lives in the Hessian of the $2$-field potential.
Stability itself is not decidable from zero-momentum data, which do not contain the momentum-dependent fluctuation operator of the full level.

The instance also shows what goes wrong with a shortcut.
One may be tempted to freeze the auxiliary field at its classical value and count the negative directions of the $\Psi \Psi$ block~\eqref{p:eq:sol:H-psi-psi} alone.
That block keeps only the direct channel, one third of the quartic term in~\eqref{p:eq:sol:channels}, and misses the exchange, which is twice as large and not suppressed at this level.
For the zero-momentum data of the companion paper~\cite{Erbin:2026:SL2CVerticesClosed} it gives the wrong count at some critical points.
At the perturbative vacuum and at the small-amplitude points, where the quartic channels do not compete with the kinetic and cubic terms, the count is right.
The small-amplitude points lie in the dilaton valley that the truncation itself creates.

The correct count is~\eqref{p:eq:sol:inertia}: the inertia of the joint Hessian minus the constant inertia of the auxiliary block.
Equivalently, freezing the fluctuation $\sigma_2^{(2)} = \Sigma_2 - \Psi^2$ instead of $\Sigma_2$ gives the right count, since in the completed-square form of \cref{p:sec:hs:quartic} the action splits exactly into the eliminated action plus a quadratic form in $\sigma_2^{(2)}$.
This is the usual distinction between a mean field and the random-phase approximation in a Hubbard--Stratonovich treatment~\cite{Negele:1988:QuantumManyParticleSystems}.
A fluctuation analysis around a mean-field saddle must therefore integrate $\sigma_2$ jointly, or use the Schur complement.

\subsubsection{Elimination of \texorpdfstring{$\Psi$}{Psi}}
\label{p:sec:sol:psi-elimination}

The cubic action is quadratic in $\Sigma_2$, which is why we eliminate $\Sigma_2$.
It is cubic in $\Psi$, so $\Psi$ cannot be eliminated in closed form, but we can ask what the equation for $\Psi$ at fixed $\Sigma_2$ looks like, since it is the equation of a Hubbard--Stratonovich mean field.
It gives a second route to the classical solutions, with the auxiliary field as the variable, and it tells whether the elimination of $\Sigma_2$ misses any of them (end of this subsection).
At quartic order it reads:
\begin{equation}
  \label{p:eq:sol:psi-equation}
  Q_{\Sigma_2} \Psi
    + \frac{\gs}{2} \, \ell_2(\Psi^2)
    = 0,
  \qquad
  Q_{\Sigma_2}
    := Q_B + \frac{\gs^2}{6} \, \ell_3(\Sigma_2, \cdot),
\end{equation}
where the terms linear in $\Psi$ define a kinetic operator deformed by the background $\Sigma_2$.
This operator is symmetric with respect to the pairing, since $\vertex{4}$ is, but it is not nilpotent for a generic $\Sigma_2$.
Its square is controlled by the $\Loo$ relation at arity three.
The $O(\gs^2)$ term is the failure of $\ell_3(\Sigma_2, \cdot)$ to commute with $Q_B$.
By the arity-three relation it is the sum of $\ell_3$ applied to $\hat Q_B \Sigma_2$, with $\hat Q_B$ the $2$-string BRST operator of \cref{p:sec:gauge}, and of the Jacobiator of $\ell_2$ on $\Sigma_2$, and it vanishes only if this sum does.
The $O(\gs^4)$ term $\ell_3(\Sigma_2, \ell_3(\Sigma_2, \cdot))$ is constrained by no relation at all.
On the solution $\Sigma_2 = \Psi^{* 2}$, $Q_{\Sigma_2}$ is not the shifted operator $Q_{\Psi^*}$ of~\eqref{p:eq:sol:channels}.
It lacks the $\ell_2$ term, being the linearization at $\Psi = 0$ rather than at $\Psi^*$, and it carries one third of the $\ell_3$ term, the direct channel again.
The quintic order adds to~\eqref{p:eq:sol:psi-equation} the term $\tfrac{\gs^3}{5!} \ell_4(\Sigma_2^2)$, a tadpole for $\Psi$ in the background $\Sigma_2$, and leaves $Q_{\Sigma_2}$ unchanged.

In Siegel gauge, the solutions of~\eqref{p:eq:sol:psi-equation} come in branches.
$\Psi = 0$ solves it for every $\Sigma_2$, and this trivial branch is isolated wherever $Q_{\Sigma_2}$ is invertible, by the implicit function theorem.
A nontrivial branch can pass through $(0, \Sigma_2)$ only where $Q_{\Sigma_2}$ has a kernel, which happens on a hypersurface in the auxiliary space.
It does pass when the kernel is one-dimensional and the cubic vertex does not vanish on it, see below.
Since every term of the action carries one power of $\gs$ per leg beyond two, with $\Sigma_2$ counting two legs, that hypersurface sits at $\Sigma_2 = O(\gs^{-2})$: the scale of the tachyon vacuum, where nothing is perturbative.
Below that scale $Q_{\Sigma_2}$ is invertible by a Neumann series, at generic momentum in Siegel gauge.

At zero momentum the massless states lie in the kernel of $Q_B$, so $\Sigma_2 = 0$ itself is on the hypersurface and the small-$\Sigma_2$ statements do not apply.
The statement that follows is local and holds at any point of the hypersurface where the kernel is one-dimensional, in either regime.
At a point of the hypersurface where the kernel is one-dimensional, spanned by $\kappa$ with $\vertex{3}(\kappa^3) \neq 0$, the crossing is transcritical.
To first order in the displacement $\delta \Sigma_2$ from the hypersurface, the nontrivial branch is:
\begin{equation}
  \label{p:eq:sol:transcritical}
  \Psi^*(\Sigma_2)
    = - \frac{\gs}{3} \, \frac{\vertex{4}(\delta \Sigma_2, \kappa^2)}{\vertex{3}(\kappa^3)} \, \kappa.
\end{equation}

Along a branch $\Psi^*(\Sigma_2)$ the action becomes a function of $\Sigma_2$ alone, the tree-level effective action $\Gamma$ of the auxiliary field:
\begin{equation}
  \label{p:eq:sol:gamma}
  \Gamma[\Sigma_2]
    := \caction \big( \Psi^*(\Sigma_2), \Sigma_2 \big),
  \qquad
  \dd \Gamma
    = \partial_{\Sigma_2} \caction \big|_{\Psi^*(\Sigma_2)},
\end{equation}
the second identity because $\partial_\Psi \caction$ vanishes on the branch, the mirror of the envelope identity used for the eliminated action.
On the trivial branch $\Gamma$ is just the quadratic term $- \tfrac{\gs^2}{4!} \vertex{4}(\Sigma_2^2)$, stationary at $\Sigma_2 = 0$ only: at quartic order there is no tree diagram with external $\Sigma_2$ legs only, since a $\Psi$ line has nowhere to end.
At quintic order the tadpole provides the end, and $\Gamma$ on the perturbative branch is the sum of tree diagrams with $\Sigma_2$ legs, the first being two tadpoles joined by a $\Psi$ propagator $- Q_{\Sigma_2}^{-1}$ at $O(\gs^6)$.
The identity in~\eqref{p:eq:sol:gamma} says that a stationary point of $\Gamma$ is a joint critical point of the cubic action.
Conversely, a joint critical point whose $\Psi$-Hessian $Q_{\Sigma_2} + \gs \, \ell_2(\Psi^*, \cdot)$ is invertible lies on a unique local branch and is stationary for its $\Gamma$.
At quartic order the joint critical points are those of the original theory, $(\Psi^*, \Psi^{* 2})$ with $\Psi^*$ a critical point of $S_{\le 4}$, so the two eliminations see the same critical points on their common domain.
Each elimination can miss only the critical points where its own Hessian block degenerates.
The elimination of $\Sigma_2$ can miss those where the dressed kinetic operator $\kin{2}{\Psi}{3}$ degenerates, from quintic order on, and the elimination of $\Psi$ those where the $\Psi$-Hessian at the point degenerates.
Neither has stationary points the other lacks, away from these two loci.

\subsubsection{The degenerate branch}
\label{p:sec:sol:degenerate}


The dressed kinetic operator of the auxiliary field depends on $\Psi$ from quintic order on, and it can fail to be invertible at some backgrounds.
At quintic order the auxiliary equation is linear, $\kin{2}{\Psi}{3}(\Sigma_2) = \sprod{2}{2}(\Psi^2)$, and its solution set $F_\Psi$ at fixed $\Psi$ is:
\begin{equation}
  \label{p:eq:sol:degenerate}
  F_\Psi
    = \begin{cases}
        \big\{ \Sigma_2^{(3)*}(\Psi) \big\}
          & \kin{2}{\Psi}{3} \text{ invertible},
        \\
        \emptyset
          & \kin{2}{\Psi}{3} \text{ singular},
            \ \sprod{2}{2}(\Psi^2) \notin \operatorname{im} \kin{2}{\Psi}{3},
        \\
        \Sigma_{2, p} + \ker \kin{2}{\Psi}{3}
          & \kin{2}{\Psi}{3} \text{ singular},
            \ \sprod{2}{2}(\Psi^2) \in \operatorname{im} \kin{2}{\Psi}{3},
      \end{cases}
\end{equation}
with $\Sigma_{2, p}$ a particular solution.
The third case is the degenerate branch: an affine space of solutions of the auxiliary equation over each $\Psi$ of the locus where the dressed operator is singular and the source lies in its image.
For a kernel of dimension one this is a line, and the two conditions cut a locus of expected codimension two in $\Psi$-space (checked in the $2$-field example with random couplings below), at $\Psi = O(1 / \gs)$ by the same scaling as above.
The branch then has one dimension less than $\Psi$-space.
The action is constant along each line, since every point of it solves the auxiliary equation, and the auxiliary field is not algebraically eliminable there, the line being a flat direction of the auxiliary block.
The regular solution $\Sigma_2^{(3)*}(\Psi)$, a rational function of $\Psi$ with the determinant of the dressed operator of the auxiliary field in the denominator, has a point of indeterminacy at each such $\Psi$ rather than a pole.
Along every direction not tangent to the singular locus of the dressed operator of the auxiliary field, the eliminated action tends there to its value on the line, which is constant.
Approaching the point along different directions, the regular solution $\Sigma_2^{(3)*}(\Psi)$ tends to different points of the line, and sweeps the whole of it: the branch is the closure of the regular solution over its base points, not a disconnected family.

It solves the auxiliary equation only.
A joint critical point on the branch needs, in addition, the $\Psi$-equation, which has as many components as $\Psi$ while the branch has one dimension less: one condition too many.

In a finite-dimensional model with generic couplings, a one-dimensional kernel and a locus of codimension two, the counting is a theorem.
It does not apply to the string, whose kinetic term is fixed, but it shows what a critical point on the branch requires: a coincidence between the kinetic term and the vertices, which the string vertices would have to produce.
We fix the cubic, quartic and quintic vertices.
The auxiliary equation does not contain the kinetic term, so the locus, the particular solution and the kernel are fixed with them, while the kinetic operator of $\Psi$ enters the $\Psi$-equation only through its value on $\Psi$.
The branch avoids $\Psi = 0$, since there the dressed operator of the auxiliary field is the quartic form, which is invertible.
For $\Psi \neq 0$ every vector is the image of $\Psi$ under some symmetric kinetic operator.
At each point of the branch, the kinetic operators that solve the $\Psi$-equation there therefore form an affine subspace whose codimension is the dimension of $\Psi$-space.
The branch has one dimension less than $\Psi$-space, so these subspaces fill a set of dimension at most one less than that of the space of kinetic operators.
That set lies in a proper algebraic subset, and for a kinetic operator outside it no joint critical point lies on the branch.

For a kinetic operator in that subset, a joint critical point on the branch, when it is nondegenerate, persists off the subset as an ordinary critical point at which the dressed operator of the auxiliary field is invertible.
It is a regular critical point crossing the locus, not a solution of its own.
In a $2$-field example with random couplings the subset has codimension one and the critical point on it is nondegenerate.

The exclusion uses the genericity of the kinetic term, and that of the quartic and quintic vertices through the two conditions on the locus; the string provides neither, its kinetic term being fixed by the spectrum.
The structure~\eqref{p:eq:sol:degenerate} and the closure statement transfer to the string vertices, under the standing assumptions.
We have not done the computation with the quintic data that would decide whether the actual vertices put the theory on that subset.
Nor have we checked whether a selection rule makes the source orthogonal to the kernel identically on the singular locus, which would change the counting.

At the level-$(1, 1)$ instance above, which stops at quartic order, the kinetic operator of the auxiliary field is the constant form $\vertex{4}$ on $\mathcal U^{\wedge 2}$, invertible, and no degenerate branch exists.
At sextic order the auxiliary equation is quadratic, and the fold described with the Hessian is the analogous locus; beyond the one-field model we have not examined it.

\subsubsection{Mean-field approximation}
\label{p:sec:sol:mean-field}


This remark is formal: it leaves the classical theory, through the path integral and an approximation that no small parameter controls.
At one loop the effective potential of the string field involves the loop vertices of the original theory, which the classical construction does not contain; without them the mean field is not a controlled approximation.
With the weight $\e^{S}$ of \cref{p:app:conventions}, the exact equation of motion of the auxiliary field in the quantum theory at quartic order is:
\begin{equation}
  \label{p:eq:sol:quantum-auxiliary}
  \sprod{2}{2}\big( \Sigma_2 - \langle \Psi^2 \rangle_{\Sigma_2} \big)
    = 0,
\end{equation}
so that the auxiliary field is, formally, the pair correlator of the string field in its own background.
We evaluate it in the Gaussian approximation, in Siegel gauge, with the cubic vertex dropped and the trace taken over the fields of ghost number two.
This gives the one-loop effective potential for $\Sigma_2$ at $\Psi = 0$ and the gap equation that follows from it:
\begin{equation}
  \label{p:eq:sol:gap}
  V[\Sigma_2]
    = \frac{\gs^2}{24} \, \vertex{4}(\Sigma_2^2)
      + \frac{\hbar}{2} \, \tr \log \big( - Q_{\Sigma_2} \big),
  \qquad
  \Sigma_2
    = \hbar \, G_{\Sigma_2}
  \quad \text{on } \Hn{2} / \ker \sprod{2}{2}.
\end{equation}
Here $G_{\Sigma_2}$ is the propagator of $Q_{\Sigma_2}$ read as a $2$-string state.
The mean field is thus a self-consistent propagator dressed by the auxiliary field, as in the Hubbard--Stratonovich derivation of the BCS gap equation~\cite{Weinberg:2005:QuantumTheoryFields-2,Negele:1988:QuantumManyParticleSystems}.
The analogy stops there: the dressing $\tfrac{\hbar \gs^2}{6} \ell_3(G, \cdot)$ is one third of the one-loop self-energy $\tfrac{\hbar \gs^2}{2} \ell_3(G, \cdot)$ of the quartic theory, Hartree and Fock terms together, and the third it keeps is the direct channel of~\eqref{p:eq:sol:channels}, the Hartree term.
The other two thirds are the diagrams with an internal $\Sigma_2$ line, which a mean field drops, and no small parameter, neither a Fermi surface nor a large $N$, suppresses them.
There is no density-of-states enhancement of the loop either, and the cubic vertex makes the slice $\Psi = 0$ non-stationary at one loop.
These obstacles to a BCS-type condensate are arguments, not theorems.
In the level-$(1, 1)$ instance, at zero momentum and with the cubic vertex dropped, the gap equation can be solved completely.
No root is a minimum of the potential on the real slice, for every value of $\hbar \gs^2$ examined.
Among its real roots, the only local convergent saddles on an admissible contour of the rotated $\Sigma_2$ integral, at the same values, are the one-loop displacements of the classical point $\Sigma_2 = \Psi^{* 2}$.
This point is the common image of the two critical points $\pm \Psi^*$ of the potential without cubic vertex.
A mean-field treatment gives back the classical structure and no new vacuum, and we expect this conclusion to persist when the loop momentum is restored.

\subsubsection{Physical meaning}
\label{p:sec:sol:meaning}

At quartic order the cubic description changes nothing about the classical solutions: the joint critical points are $(\Psi^*, \Psi^{* 2})$ with $\Psi^*$ a critical point of $S_{\le 4}$.
At higher orders they are the critical points of the eliminated action, which agrees with the original one through the order kept, as a statement about Taylor coefficients in $\gs$.
At quintic order, in finite-dimensional models with a generic kinetic term, a one-dimensional kernel and a singular locus of codimension two, the only additional configurations, the lines of the degenerate branch, solve the auxiliary equation alone and carry no critical point.
Whether the string vertices place a critical point on them is open, and the sextic order is not covered.

The cubic description changes the fluctuation problem.
Around any solution, the second variation is one matrix, affine in the background and built from the couplings of the cubic action, with no inverse taken around the background.
Through quintic order these couplings are the elementary vertices, with no inverse of $\sprod{2}{2}$ and no reduced vertex.
The second variation is congruent to the Hessian of the eliminated action, which is that of the original theory through the order kept, plus an algebraic auxiliary block, constant at quartic order and linear in the background from quintic order on.
Wherever that block is invertible, the counts of negative and zero directions of the cubic theory are those of the eliminated theory plus those of the auxiliary block, and the mass shells are those of the eliminated theory.
The auxiliary block is the price of the cubic form: an indefinite sector whose quadratic term contains no $L_0^+$, whose negative directions are rotated in the path integral and subtracted in every count.
Whether it has mass shells at nonzero momentum is open.
Freezing the auxiliary field at its classical value keeps only the direct channel and can give the wrong count of negative directions.
The singularities of the eliminated description, the fold and the points of indeterminacy of the auxiliary solution, are artifacts of the elimination.


\subsection{The dilaton direction of the level-\texorpdfstring{$(1, 1)$}{(1,1)} truncation at quintic order}
\label{sec:sol:dilaton-d5}

\status{Conditional(PH27, SL: PH5) -- mixed}\record{095-d5-five-field-truncation (R89)}

In the level-$(1, 1)$ instance the dilaton is not flat.
Along the dilaton direction the quartic truncation leaves a nonzero term of order $d^4$ in the potential, which in the full theory is cancelled by the exchange of a massive state the truncation drops, by the dilaton identity of the companion paper at that order.
We ask the same question one order higher: in the five-field truncation of the quintic cubic action, is the coefficient of $d^5$ zero, and does the auxiliary field contribute to it?
The answer, under the assumptions below, is that the coefficient is nonzero, that it is again the exchange of the state the truncation drops, and that the auxiliary field contributes nothing below degree six in $\Psi$.

\paragraph{Assumptions and scope.}

The computation is at tree level and zero momentum, and the coefficients are formal Taylor coefficients in $d$.
The field content is that of the level-$(1, 1)$ instance, $\Psi = t \, T + d \, D$ and $\Sigma_2 \in \mathcal U^{\wedge 2}$ with $\mathcal U = \Span(T, D)$, and the action is the quintic cubic action $\caction[\le 5]$ of \cref{p:sec:hs:quintic} restricted to it, with the quartic form invertible on $\mathcal U^{\wedge 2}$.
The algebra of the first three steps below holds for arbitrary couplings.
The identification of the coefficient and its non-vanishing use two further assumptions.
\begin{itemize}
  \setlength{\itemsep}{0.5ex}

  \item The vertices kept, the quintic one included, are the truncation of a consistent genus-zero family of string vertices.
    This includes the normalization of the quintic vertex relative to the cubic and quartic ones, fixed through factorization on the boundary of the moduli space, $C_5 = C_4^2 / C_3$ for the measure constants in each normalization.
    It follows that the quintic vertex converts from the normalization of the companion paper to ours by the same factor $\tfrac{1}{2}$ per sphere correlator as the lower vertices, for which this factor is established.
    The assumption is open: no consistent quintic vertex of the family has been constructed.

  \item The order-$d^5$ dilaton identity of the companion paper~\cite{Erbin:2026:SL2CVerticesClosed}, which rests on that paper's assumption that its five-point vertex extends the family consistently, and which it establishes numerically.
\end{itemize}
The section of the vertex family and the two values of the family parameter are those of the companion paper.

\paragraph{The couplings.}

We write
\begin{equation}
  \label{eq:sol:d5:couplings}
  m
    := \pair{T}{Q_B T},
  \qquad
  b
    := \vertex{3}(T, D^2),
  \qquad
  c
    := \vertex{3}(T^2, D),
  \qquad
  h_k
    := \vertex{4}(T^{4 - k}, D^k),
  \qquad
  f_k
    := \vertex{5}(T^{5 - k}, D^k).
\end{equation}
In our conventions $m = 1$.
The kinetic term does not mix the two fields and vanishes on the dilaton, $\pair{T}{Q_B D} = \pair{D}{Q_B D} = 0$, and the cubic couplings $\vertex{3}(T^2, D)$ and $\vertex{3}(D^3)$ vanish by the chiral ghost number.
We keep $c$ symbolic in the derivation, to see which term it would produce, and set it to zero at the end.
On $\mathcal U^{\wedge 2}$ we use the basis
\begin{equation}
  e_1
    = T \otimes T,
  \qquad
  e_2
    = \tfrac{1}{2} (T \otimes D + D \otimes T),
  \qquad
  e_3
    = D \otimes D,
\end{equation}
in which $\Sigma_2 = \sum_i \sigma_i \, e_i$ and $\Psi^2$ has the components $p = (t^2, 2 t d, d^2)$, as in~\eqref{p:eq:sol:level11}.
The quartic form and the quintic vertex with one leg on $\Psi$ become the symmetric matrices
\begin{equation}
  A_{ij}
    := \vertex{4}(e_i, e_j),
  \qquad
  W(\Psi)_{ij}
    := \vertex{5}(e_i, e_j, \Psi)
    = t \, f_{i + j - 2} + d \, f_{i + j - 1},
\end{equation}
the second equality because the legs of $e_i$ carry $i - 1$ dilatons.

\paragraph{Step 1: the auxiliary field adds nothing below degree six.}

With these notations the quintic cubic action~\eqref{p:eq:hs5:S5tilde} restricted to the truncation reads
\begin{equation}
  \label{eq:sol:d5:action}
  \caction[\le 5]
    = \frac{m}{2} \, t^2
      + \frac{\gs}{6} \big( a \, t^3 + 3 c \, t^2 d + 3 b \, t d^2 \big)
      - \frac{\gs^2}{24} \, \sigma^\top A \sigma
      + \frac{\gs^2}{12} \, \sigma^\top A p
      + \frac{\gs^3}{120} \, \sigma^\top W(\Psi) \sigma,
\end{equation}
with $a := \vertex{3}(T^3)$, which will not enter.
The terms quadratic in $\sigma$ combine into $- \tfrac{\gs^2}{24} \, \sigma^\top K \sigma$ with the dressed operator
\begin{equation}
  K(\Psi)
    := A - \frac{\gs}{5} \, W(\Psi),
\end{equation}
the truncated form of the operator of $\Sigma_2$ at quintic order.
The auxiliary equation $K \sigma = A p$ gives $\sigma^* = K^{-1} A p$, and on this solution
\begin{equation}
  \label{eq:sol:d5:sigma-eliminated}
  \begin{aligned}
    \caction[\le 5] \big|_{\sigma = \sigma^*}
      &
      = S_{\le 3}(\Psi)
        + \frac{\gs^2}{24} \, p^\top A K^{-1} A p
    \\
      &
      = S_{\le 3}(\Psi)
        + \frac{\gs^2}{24} \bigg[
          p^\top A p
          + \sum_{k \ge 1} \Big( \frac{\gs}{5} \Big)^k \, p^\top W (A^{-1} W)^{k - 1} p
        \bigg]
    \\
      &
      = S_{\le 5} \big|_{\mathcal U}
        + \frac{\gs^4}{600} \, v^\top A^{-1} v
        + O(\gs^5 \Psi^7),
      \qquad
      v
        := W(\Psi) \, p,
  \end{aligned}
\end{equation}
where $S_{\le 3}$ is the kinetic term plus the cubic vertex.
The second line expands $K^{-1}$ in its Neumann series around $A^{-1}$.
Since $W$ is linear in $\Psi$ and $p$ is quadratic, the $k$-th term has degree $4 + k$ in $\Psi$.
The term $p^\top A p$ is $\vertex{4}(\Psi^4)$ and the term $k = 1$ is $p^\top W p = \vertex{5}(\Psi^5)$, which give the quartic and quintic vertices of the original action with their coefficients $\gs^2 / 4!$ and $\gs^3 / 5!$.
The first term of the cubic formulation of its own is the one of degree six, the composite term of~\eqref{p:eq:sol:matching}.
As a statement about Taylor coefficients this is exact; the series itself converges for $\gs \lVert A^{-1} W(\Psi) \rVert < 5$.
The auxiliary field therefore contributes nothing to the coefficient of $d^5$ beyond the quintic vertex $\vertex{5}(D^5)$ that it reproduces.

\paragraph{Step 2: which terms reach $d^5$.}

We now eliminate $t$ along the dilaton.
Since the kinetic term does not mix $T$ and $D$, the tachyon has no source linear in $d$, and its solution starts at $t^* = O(\gs d^2)$.
A monomial $t^j d^n$ then contributes from order $d^{n + 2 j}$ on.
The monomials of weight at most five are $t^2$, which enters only through the propagator $1 / m$, and $t d^2$, $t d^3$, $t^2 d$, $d^4$, $d^5$.
The monomials $t^3$, $t^2 d^2$ and the whole degree-six remainder of~\eqref{eq:sol:d5:sigma-eliminated} have weight six or more.
This counting uses $\pair{T}{Q_B D} = \pair{D}{Q_B D} = 0$ and $\vertex{3}(D^3) = 0$: without them, $t$ would be sourced at order $d$, or $d$ would not be flat at quadratic and cubic order.
Three structures can thus reach $d^5$: the contact term $\vertex{5}(D^5)$, the exchange of one tachyon between $\vertex{3}(T, D^2)$ and $\vertex{4}(T, D^3)$, and a term with two tachyons through $\vertex{3}(T^2, D)$.

\paragraph{Step 3: the coefficient.}

The $t$-equation reads
\begin{equation}
  \partial_t \caction
    = m \, t
      + J(d)
      + \gs c \, t d
      + O(t^2, \gs^2 t d^2),
  \qquad
  J(d)
    := \frac{\gs}{2} \, b \, d^2
      + \frac{\gs^2}{6} \, h_3 \, d^3,
\end{equation}
so that $t^* = - \tfrac{\gs b}{2 m} \, d^2 + O(d^3)$.
The correction of order $d^3$ to $t^*$ does not enter at order $d^5$: it multiplies $\partial_t \caction$ at order $d^2$, which vanishes on the solution (the $2 n + 1$ rule).
Inserting $t^*$ gives, through $d^5$,
\begin{equation}
  \caction[\rm eff](d)
    = \frac{\gs^2}{24} \, h_4 \, d^4
      + \frac{\gs^3}{120} \, f_5 \, d^5
      - \frac{J(d)^2}{2 m}
      + \frac{\gs}{2} \, c \, t^{* 2} d
      + O(d^6),
\end{equation}
where $- J^2 / 2 m$ is $\tfrac{m}{2} t^{* 2} + t^* J$ at the orders kept.
Its expansion is $- \tfrac{\gs^2 b^2}{8 m} \, d^4 - \tfrac{\gs^3 b \, h_3}{12 m} \, d^5$, and the last term is $\tfrac{\gs^3 c \, b^2}{8 m^2} \, d^5$.
On the potential side, $V_{\rm eff} = - \caction[\rm eff]$, the coefficients of $d^4$ and $d^5$ are
\begin{equation}
  \label{eq:sol:d5:c5}
  V_{\rm eff}(d)
    \supset \gs^2 \Big( \frac{b^2}{8 m} - \frac{h_4}{24} \Big) d^4
      + \gs^3 c_5 \, d^5,
  \qquad
  c_5
    = - \frac{f_5}{120}
      + \frac{b \, h_3}{12 m}
      - \frac{c \, b^2}{8 m^2}.
\end{equation}
With $c = \vertex{3}(T^2, D) = 0$ the term with two tachyons drops, and
\begin{equation}
  \label{eq:sol:d5:c5-final}
  c_5
    = - \frac{\vertex{5}(D^5)}{120}
      + \frac{\vertex{3}(T, D^2) \, \vertex{4}(T, D^3)}{12 \, \pair{T}{Q_B T}}.
\end{equation}
The couplings $\vertex{3}(T^3)$, $\vertex{4}(T^4)$, $\vertex{4}(T^3, D)$, $\vertex{4}(T^2, D^2)$ and $\vertex{5}(T^{5 - k}, D^k)$ for $k \le 4$ do not enter.
This closed form was checked by eliminating $\Sigma_2$ and $t$ exactly on random rational couplings, with $c \neq 0$ and $m \neq 1$, which also confirms that the remainder of~\eqref{eq:sol:d5:sigma-eliminated} has no term of degree five or less.
The $d^4$ coefficient is the one of the quartic instance.

\paragraph{Step 4: the same coefficient as a channel count.}

The coefficient can be compared with the dilaton identity of the companion paper once it is written in that paper's normalization.
There the couplings are written on the potential side, as brackets $\{X\}$, and each vertex is twice ours, so that $\vertex{n}(X) = - \tfrac{1}{2} \{X\}$; in particular $m = - \tfrac{1}{2} \{T T\} = 1$.
Substituting into~\eqref{eq:sol:d5:c5},
\begin{equation}
  \label{eq:sol:d5:brackets}
  \begin{aligned}
    c_5
      &
      = \frac{1}{2} \cdot \frac{1}{120} \Big[
        \{D^5\}^{\rm contact}
        + 10 \, \Pi_T \, \{D D T\} \{T D^3\}
        + 15 \, \Pi_T^2 \, \{D D T\}^2 \{T T D\}
      \Big],
      \qquad
      \Pi_T
        = - \{T T\}^{-1}
        = \frac{1}{2},
    \\
    240 \, c_5
      &
      = \{D^5\}^{\rm contact}
        + 5 \, \{D D T\} \{T D^3\}
        + \frac{15}{4} \, \{D D T\}^2 \{T T D\},
      \qquad
      \{T T D\}
        = 0.
  \end{aligned}
\end{equation}
The coefficients $10$ and $15$ are the channel counts of the tree diagrams with five external dilatons: $\binom{5}{2} = 10$ pairs of legs on the cubic vertex for one tachyon propagator, and $15$ trees with two propagators.
This form is the same elimination of $t$ written as a sum over channels, not an independent derivation; its agreement with~\eqref{eq:sol:d5:c5} restates Step 1 together with the sign conventions.
What constrains the factor $5$ is the coefficient $\gs^3 / 5!$ of the quintic term in the cubic action: with $\gs^3 / 4!$ in its place the elimination no longer reproduces the channel count.
The sign of the exchange is fixed by $\Pi_T > 0$, that is, by the tachyonic $T$, with $\pair{T}{Q_B T} = + 1$ in the action and $V = - S$.
Every structure of the truncation has its counterpart, with the same coefficient, in the order-$d^5$ dilaton identity of the companion paper, which reads $\{D^5\}^{\rm contact} + 5 \, \{D D T\} \{T D^3\} - 5 \, \{D D f_1\} \{f_1 D^3\} = 0$.
The one term of that identity which~\eqref{eq:sol:d5:brackets} does not contain is the exchange of the level-$(2, 2)$ state
\begin{equation}
  f_1
    = c_{-1} \bar c_{-1} \ket{0},
\end{equation}
which enters with the opposite sign since its propagator is $\Pi_{f_1} = - \tfrac{1}{2}$.

\paragraph{Result.}

Under the two assumptions above, the identity gives
\begin{equation}
  \label{eq:sol:d5:result}
  c_5
    = \frac{\{D^5\}^{\rm contact} + 5 \, \{D D T\} \{T D^3\}}{240}
    = \frac{\{D D f_1\} \{f_1 D^3\}}{48}.
\end{equation}
With the values of the companion paper, which we do not reproduce, the two sides of~\eqref{eq:sol:d5:result} agree to the precision of that paper's own identity, and $c_5$ is positive at both values of the family parameter.
The two pieces the truncation keeps, the contact term and the tachyon exchange, nearly cancel, and what remains is the exchange of $f_1$.
The truncated dilaton valley, lifted at order $d^4$, is therefore tilted at order $d^5$.
The two terms are of comparable size only at $\abs{\gs d} \sim 1$, outside the regime of the expansion, so the stationary point they would produce together says nothing about the truncation.

The agreement of the factor $5$ and of the signs is a consistency check, under the factorization rule of the first assumption, that the coefficients $2 \gs^2 / 4!$ and $\gs^3 / 5!$ of the cubic action agree with the Feynman rules and the sign conventions.
It does not test the normalization of the quintic vertex: every term of the identity carries the weight of one sphere correlator, so that, given the identity in the companion paper's normalization, the identity in ours is equivalent to the factorization rule, and nothing here fixes the conversion factor of the quintic vertex independently of that rule.

\paragraph{Scope of the exclusion.}

The negative statement is that the five-field truncation has no flat dilaton direction at order $d^5$.
The class is the five-field truncation at zero momentum and tree level, in the section and at the two values of the family parameter of the companion paper; it is the lowest level that contains the dilaton.
Outside it, the six-field truncation $\{T, D, f_1\}$ and every truncation to higher levels are flat through $d^5$, by the results of the companion paper, at $d^5$ under its assumption on the five-point vertex.
On the level-$(2, 2)$ Siegel-gauge states the coupling to two dilatons is proportional to the pairing with $f_1$ alone, so the full exchange at that level is the exchange of $f_1$.
In the six-field truncation $\{T, D, f_1\}$ the weight counting of Step 2 is unchanged, $f_1$ being sourced at order $d^2$; the couplings $\{T D f_1\}$ and $\{f_1 D f_1\}$ vanish by the cubic chiral ghost-number rule, and no state beyond level $(2, 2)$ is sourced.
The $d^4$ and $d^5$ coefficients of the six-field truncation are thus the full tree-level ones, zero at $d^4$ exactly and at $d^5$ numerically, under the companion paper's assumption.
Step 1 holds for any $\mathcal U$ with the auxiliary field in $\mathcal U^{\wedge 2}$, so the cubic formulation adds nothing at $d^5$ there either.
The six-field action itself, with its couplings in our normalization, is not computed here: the statement concerns its coefficients through $d^5$ only.
In other sections of the vertex family the value of $c_5$ changes, and its non-vanishing is not established there.
The weakest relaxation that evades the exclusion is to add the single field $f_1$.

\paragraph{Physical meaning.}

The flatness of the dilaton at tree level is a property of the full theory, in which the contact couplings of the dilaton are balanced by the exchange of the other states: at the orders computed here, the tachyon and the massive state $f_1$ (at $d^5$ under the assumption on the five-point vertex stated above).
A level truncation cuts the balance.
At level $(1, 1)$ the only state missing from it, at orders $d^4$ and $d^5$ alike, is $f_1$, so the defect of the truncation along the dilaton is the exchange of one state and the remedy is one field.
The auxiliary field cannot supply the missing exchange: it lives in the symmetric square of the same truncated space, and its first contribution of its own has degree six.
This is the quintic-order counterpart of the statement of the quartic instance, that nothing specific to the cubic formulation enters before the quintic vertex: at quintic order nothing enters before degree six.

\paragraph{The truncation lemma.}

The last block is the appendix of the earlier version of the paper that states and proves the truncation lemma used in the subsection on level truncations, in the same notation.
The universality assumption it states is the rule discussed before the section on classical solutions, where its status is given.

\subsection{Proof of the truncation lemma}
\label{p:app:truncation}


In this appendix, we state and prove the lemma used in \cref{p:sec:solutions} to relate the solutions of a level truncation of the cubic action to those of the full cubic action and of the original theory.
The argument is algebraic: we treat the vertices as given multilinear forms, and their worldsheet definition enters through one assumption, stated below, in addition to the standing assumptions on the quartic form.

\subparagraph{Setting.}

Let $\mathcal U \subset \Hs$ be a finite-dimensional subspace of zero-momentum states of ghost number two, and let $\mathcal U^c$ be a complement:
\begin{equation}
  \label{p:eq:trunc:split-H}
  \Hs
    = \mathcal U \oplus \mathcal U^c.
\end{equation}
In the application, $\mathcal U$ is a level truncation: it lies in the zero-momentum, twist-even universal subspace and is spanned by its Siegel-gauge states of a finite set of levels, such as the tachyon and the ghost dilaton.
Its complement $\mathcal U^c$ contains the universal states outside $\mathcal U$ and the non-universal states, among them all states of nonzero momentum.
We assume that the kinetic term does not mix $\mathcal U$ with $\mathcal U^c$, as in~\eqref{p:eq:sol:kinetic-block}.
The $2$-string space then splits according to the number of strands lying outside $\mathcal U$:
\begin{equation}
  \label{p:eq:trunc:split-H2}
  \Hn{2}
    = \mathcal U^{\wedge 2}
      \oplus \big( \mathcal U \wedge \mathcal U^c \big)
      \oplus \big( \mathcal U^c \big)^{\wedge 2},
\end{equation}
where $\mathcal U^{\wedge 2}$ is spanned by the symmetrized products of two states of $\mathcal U$ and $(\mathcal U^c)^{\wedge 2}$ by those of two states of $\mathcal U^c$.
The middle summand is spanned by the symmetrized products of one state of each:
\begin{equation}
\label{p:eq:trunc:wedge}
\mathcal U \wedge \mathcal U^c
  := \Span \{ u \otimes v + v \otimes u \,:\, u \in \mathcal U, \ v \in \mathcal U^c \}.
\end{equation}

A level truncation of the cubic action restricts the string field to $\mathcal U$ and extends to the auxiliary field by restricting it to $\mathcal U^{\wedge 2}$, and the truncated action is the cubic action evaluated on these fields:
\begin{equation}
  \label{p:eq:trunc:action}
  \caction[\le n]^{[\mathcal U]}(\Psi, \Sigma_2)
    := \caction[\le n](\Psi, \Sigma_2),
  \qquad
  \Psi \in \mathcal U,
  \quad
  \Sigma_2 \in \mathcal U^{\wedge 2},
\end{equation}
for $n = 4, 5, 6$, where the cubic action at sextic order is:
\begin{equation}
  \label{p:eq:trunc:action-explicit}
  \begin{aligned}
    \caction[\le 6]
      = \frac{1}{2} \, \pair{\Psi}{Q_B \Psi}
        &
        + \frac{\gs}{3!} \, \vertex{3}(\Psi^3)
        - \frac{\gs^2}{4!} \, \vertex{4}(\Sigma_2^2)
        + \frac{2 \gs^2}{4!} \, \vertex{4}(\Sigma_2, \Psi^2)
      \\
      &
      + \frac{\gs^3}{5!} \, \vertex{5}(\Sigma_2^2, \Psi)
      + \frac{\gs^4}{6!} \, \dvertex{6}(\Sigma_2^3),
  \end{aligned}
\end{equation}
and $\caction[\le 5]$, $\caction[\le 4]$ obtained by dropping the last term, then the last two.
The sextic term is written with $\Sigma_2$ as the only auxiliary field, as a cubic coupling of three $\Sigma_2$ through the reduced vertex of \cref{p:sec:hs:sextic}:
\begin{equation}
\label{p:eq:trunc:W6}
\dvertex{6}
  = \vertex{6} - \frac{6}{5} \, \vertex{5 \cdot 5}.
\end{equation}
Its exchange part $\vertex{5 \cdot 5}$, the composite vertex of \cref{p:sec:hs:quintic}, contains the inverse of $\sprod{2}{2}$ on the whole of $\Hn{2}$.
Here we read $\dvertex{6}$ and $\vertex{5 \cdot 5}$ as trilinear forms on $\Hn{2}$, symmetrized in their three $2$-string arguments (\cref{p:sec:hs:quintic,p:sec:hs:sextic}).
We do not treat here the sextic action with $\Sigma_3$ as a second auxiliary field.
Since the strands of $\Sigma_2$ are Grassmann even, no sign enters anywhere in this appendix.

We write the variation of the full cubic action along $\delta \Psi = \psi \in \Hs$ and $\delta \Sigma_2 = \sigma_2 \in \Hn{2}$ as:
\begin{equation}
\label{p:eq:trunc:variation}
\delta \caction[\le 6]
  = \pair{\psi}{\mathcal E^{(4)}_\Psi}
    + \frac{\gs^2}{12} \, \pair[2]{\sigma_2}{\mathcal E^{(4)}_{\Sigma_2}},
\end{equation}
which defines the two equations of motion with the normalization:
\begin{subequations}
\label{p:eq:trunc:eom}
\begin{align}
  \label{p:eq:trunc:eom-psi}
  \pair{\psi}{\mathcal E^{(4)}_\Psi}
    &
    = \pair{\psi}{Q_B \Psi}
      + \frac{\gs}{2} \, \vertex{3}(\psi, \Psi^2)
      + \frac{\gs^2}{6} \, \vertex{4}(\psi, \Psi, \Sigma_2)
      + \frac{\gs^3}{120} \, \vertex{5}(\psi, \Sigma_2^2),
  \\
  \label{p:eq:trunc:eom-sigma}
  \pair[2]{\sigma_2}{\mathcal E^{(4)}_{\Sigma_2}}
    &
    = \vertex{4}(\sigma_2, \Psi^2 - \Sigma_2)
      + \frac{\gs}{5} \, \vertex{5}(\sigma_2, \Sigma_2, \Psi)
      + \frac{\gs^2}{20} \, \dvertex{6}(\sigma_2, \Sigma_2^2),
\end{align}
\end{subequations}
whose leading terms are those of \cref{p:sec:hs:quartic}.
The superscript is the highest power of $\gs$ in the action.
At quintic order the last term of~\eqref{p:eq:trunc:eom-sigma} is absent.
At quartic order its last two terms are absent, and so is the last term of~\eqref{p:eq:trunc:eom-psi}.
We say that the equation of $\Sigma_2$ is violated at relative order $\gs^k$ along $\sigma_2$ when $\pair[2]{\sigma_2}{\mathcal E^{(n)}_{\Sigma_2}}$, normalized as in~\eqref{p:eq:trunc:eom-sigma}, is of order $\gs^k$.
The truncated equations of motion are:
\begin{equation}
\label{p:eq:trunc:eom-truncated}
\pair{\psi}{\mathcal E^{(n)}_\Psi}
  = 0
\quad
\forall \psi \in \mathcal U,
\qquad
\pair[2]{\sigma_2}{\mathcal E^{(n)}_{\Sigma_2}}
  = 0
\quad
\forall \sigma_2 \in \mathcal U^{\wedge 2},
\end{equation}
whereas the full equations require the same for all $\psi \in \Hs$ and all $\sigma_2 \in \Hn{2}$.
We want to know which of the remaining components, along $\mathcal U^c$ and along the two other summands of~\eqref{p:eq:trunc:split-H2}, a solution of the truncated equations satisfies.
\begin{equation}
  \label{p:eq:trunc:nondegenerate}
  \vertex{4}(A, B)
    = 0
  \ \ \forall B \in \mathcal U^{\wedge 2}
  \quad \Longrightarrow \quad
  A
    = 0,
  \qquad
  A \in \mathcal U^{\wedge 2},
\end{equation}
which replaces, inside the truncation, the quotient by $\ker \sprod{2}{2}$ of the untruncated theory.
If the restricted form is degenerate, we pass to the quotient of $\mathcal U^{\wedge 2}$ by its kernel, which is consistent exactly when $\vertex{5}$ vanishes on the kernel, as in the untruncated theory.
We also use that $\Psi^2 \in \mathcal U^{\wedge 2}$ whenever $\Psi \in \mathcal U$.
This is what makes $\mathcal U^{\wedge 2}$ the minimal auxiliary space: with an auxiliary space not containing $\Psi^2$, the matching generically fails already at quartic order.

Some statements use the inverse of $\sprod{2}{2}$ on all of $\Hn{2}$: the criterion in (iii), the full reduced vertex, the sextic-order residual~\eqref{p:eq:trunc:residual-sextic} and the comparison~\eqref{p:eq:trunc:schur}--\eqref{p:eq:trunc:schur-difference}.
They rest on the standing assumptions of \cref{p:sec:hs:quartic}: the quartic form is invertible on the quotient by its kernel, level by level, and the sum over intermediate $2$-string states converges.
In the comparison we take the kernel to vanish.
The proofs are exact in any finite-dimensional truncation of $\Hs$ containing $\mathcal U$.
We do not address the limit of infinite level.

\subparagraph{Universality assumption.}

The lemma rests on the universality assumption~\eqref{p:eq:sol:universality}: a vertex vanishes whenever exactly one of its legs lies in the space $\mathcal N$ of non-universal states, those of nonzero momentum included, and all the others lie in $\mathcal U$.
The kinetic term does not couple $\mathcal U$ to $\mathcal N$ either, as noted in \cref{p:sec:solutions}.
Couplings with two or more non-universal legs are not constrained, nor are those whose single non-universal leg comes with a leg outside $\mathcal U$, nor those whose legs are all universal.
The pairs of a state of $\mathcal U$ and a non-universal state form the subspace $\mathcal U \wedge \mathcal N \subset \mathcal U \wedge \mathcal U^c$, defined as in~\eqref{p:eq:trunc:wedge}.

This property allows Sen to restrict the open string field to the universal subspace~\cite{Sen:1999:UniversalityTachyonPotential}.
This subspace is spanned by the states built on the zero-momentum vacuum by ghost oscillators and matter Virasoro generators, and on it the tachyon potential is independent of the background.
The argument is a Ward identity: a vertex is a moduli integral of a sphere correlator of conformally transformed states with $b$-ghost insertions.
The conformal maps preserve the matter Virasoro module of each leg, and the $b$-insertions act on the ghost factor only.
The matter factor is then a correlator with exactly one state outside the identity module, which vanishes.
The same argument applies to the closed string and underlies the level-truncation studies of the closed-string tachyon vacuum~\cite{Yang:2005:ClosedStringTachyon}, where the matter primaries are dropped and the ghost dilaton is kept.

The rule is established at genus zero and zero momentum, on the Fock space of zero momentum, for any vertex given as such a moduli integral, and it holds for the cubic and quartic vertices of the $\group{SL}(2, \C)$ family in any admissible section; at five points and more it rests on this family-independent Ward argument, since no quintic vertex of that family exists yet.\footnote{Ledger entry \recordref{R88}; a cited result (Yang and Zwiebach, after Sen), with no derivation record of its own.}
We use it in the form~\eqref{p:eq:sol:universality}.
It is falsified by a nonzero coupling with a single non-universal leg and all other legs in $\mathcal U$, such as $\vertex{4}(T^3, x)$ with $T$ the tachyon and $x$ a zero-momentum matter primary.
It is needed only at zero momentum: a coupling with exactly one leg at nonzero momentum and all others at zero momentum vanishes by momentum conservation.
The assumption concerns the elementary vertices $\vertex{n}$, not the reduced vertex $\dvertex{6}$, whose exchange part need not vanish with a single non-universal leg.
This exchange part is the source of the single sextic-order residual found below.

\subparagraph{Lemma.}

\emph{Let $\mathcal U$ lie in the zero-momentum universal subspace, let the quartic vertex be non-degenerate on $\mathcal U^{\wedge 2}$, and let the kinetic block~\eqref{p:eq:sol:kinetic-block} and the universality assumption~\eqref{p:eq:sol:universality} hold.
  Let $(\Psi^*, \Sigma_2^*) \in \mathcal U \times \mathcal U^{\wedge 2}$ be a solution of the truncated equations of motion~\eqref{p:eq:trunc:eom-truncated}, on the branch connected to $\Sigma_2 = \Psi^2$ at $\gs = 0$.
  Then:}
\begin{enumerate}[label=(\roman*)]
  \setlength{\itemsep}{0.5ex}

  \item \emph{the equation of motion of $\Psi$ of the full cubic action holds along every direction of $\mathcal N$, at quartic, quintic and sextic order, identically for every $(\Psi, \Sigma_2) \in \mathcal U \times \mathcal U^{\wedge 2}$ and not only on-shell.
      Along a universal direction outside $\mathcal U$ it is in general violated from order $\gs$, by the level-truncation residual~\eqref{p:eq:trunc:level-residual};}

  \item \emph{at quartic order, the equation of motion of $\Sigma_2$ of the full cubic action holds along every direction of $\Hn{2}$.
      The truncated solution thus solves every equation of motion of $\caction[\le 4]$ except that of $\Psi$ along the universal directions outside $\mathcal U$, whose residual is that of level truncation;}

  \item \emph{at quintic order, the equation of motion of $\Sigma_2$ holds identically along $\mathcal U \wedge \mathcal N$, the pairs of a state of $\mathcal U$ and a non-universal state.
      Along the other pairs outside $\mathcal U^{\wedge 2}$ the truncated equations do not imply it: its residual~\eqref{p:eq:trunc:residual} is of relative order $\gs$ and vanishes along all of them exactly when $\sprod{2}{2}^{-1} \circ \sprod{3}{2}(\Psi^3)$ lies in $\mathcal U^{\wedge 2}$.
      At sextic order the component along $\mathcal U \wedge \mathcal N$ is no longer protected: it is of relative order $\gs^2$, through the exchange part of $\dvertex{6}$, unless a symmetry removes it;}

  \item \emph{eliminating $\Sigma_2$ inside the truncation reproduces the original action restricted to $\mathcal U$ exactly at quartic order and through $O(\gs^3)$ at quintic order, as the elimination in the full theory does.
      At sextic order it does so through $O(\gs^4)$, provided the reduced vertex is built with the inverse of the quartic form on $\mathcal U^{\wedge 2}$ rather than on $\Hn{2}$.
      This truncated action is no longer the restriction of $\caction[\le 6]$, and its solution violates the equation of $\Sigma_2$ of $\caction[\le 6]$ along $\mathcal U^{\wedge 2}$ at relative order $\gs^2$.}
\end{enumerate}
\emph{The orders in (iii) and (iv) are orders of the Taylor expansion in $\gs$ at fixed $\Psi$.}

The universality assumption is used only for the directions of $\mathcal N$ in (i) and for $\mathcal U \wedge \mathcal N$ in (iii).
The other statements need only the non-degeneracy~\eqref{p:eq:trunc:nondegenerate} and the kinetic block.
Statements (iii) and (iv) use only the truncated equation of $\Sigma_2$, at any fixed $\Psi \in \mathcal U$.

Statement (i) means that the truncated solution sources no non-universal field.
Along the universal states outside $\mathcal U$, its error is that of a level truncation of the original theory.
The auxiliary field does need to leave $\mathcal U^{\wedge 2}$ from quintic order on.
At quintic order it does so, in general, along the pairs outside $\mathcal U^{\wedge 2}$ other than $\mathcal U \wedge \mathcal N$, and at sextic order also along $\mathcal U \wedge \mathcal N$, through the exchange part of $\dvertex{6}$.
The error this causes in the action is of the order at which the truncated elimination stops matching the full one.

\subparagraph{Proof.}

\emph{(i) The $\Psi$-equation along $\mathcal N$ and along the universal states outside $\mathcal U$.}
We take $\psi \in \mathcal N$ in~\eqref{p:eq:trunc:eom-psi}, with $\Psi \in \mathcal U$ and $\Sigma_2 \in \mathcal U^{\wedge 2}$.
By multilinearity, expanding $\Sigma_2 = \sum_a \lambda_a \, u_a \otimes u'_a$ with $u_a, u'_a \in \mathcal U$, every term of~\eqref{p:eq:trunc:eom-psi} becomes a sum of vertices with all legs in $\mathcal U$ except the single non-universal leg $\psi$:
\begin{equation}
\label{p:eq:trunc:psi-outside}
\begin{aligned}
  \pair{\psi}{\mathcal E^{(4)}_\Psi}
    &
    = \pair{\psi}{Q_B \Psi}
      + \frac{\gs}{2} \, \vertex{3}(\psi, \Psi, \Psi)
      + \frac{\gs^2}{6} \sum_a \lambda_a \, \vertex{4}(\psi, \Psi, u_a, u'_a)
    \\
    &
    \quad
    + \frac{\gs^3}{120} \sum_{a, b} \lambda_a \lambda_b \, \vertex{5}(\psi, u_a, u'_a, u_b, u'_b).
\end{aligned}
\end{equation}
The first term vanishes by the kinetic block~\eqref{p:eq:sol:kinetic-block}.
The others vanish by the universality assumption, term by term and for every value of the fields.
The sextic term of the action contains no $\Psi$ leg, so (i) holds at sextic order as well.

Along a universal direction $\psi$ outside $\mathcal U$ the assumption gives nothing.
The kinetic term drops by~\eqref{p:eq:sol:kinetic-block}, and at the truncated solution the residual is:
\begin{equation}
\label{p:eq:trunc:level-residual}
\pair{\psi}{\mathcal E^{(n)}_\Psi}
  = \frac{\gs}{2} \, \vertex{3}(\psi, \Psi^{* 2})
    + \frac{\gs^2}{6} \, \vertex{4}(\psi, \Psi^{* 3})
    + O(\gs^3),
\qquad
n = 2, 3, 4,
\end{equation}
which is the residual of the level-truncated original theory along $\psi$~\cite{Yang:2005:ClosedStringTachyon}, exactly at quartic order and through order $\gs^2$ beyond.
At quartic order $\Sigma_2^* = \Psi^{* 2}$ exactly, by (ii) below, and the two terms shown are the whole residual.
Beyond, $\Sigma_2^* = \Psi^{* 2} + O(\gs)$, so the term $\tfrac{\gs^2}{6} \vertex{4}(\psi, \Psi^*, \Sigma_2^*)$ equals $\tfrac{\gs^2}{6} \vertex{4}(\psi, \Psi^{* 3})$ up to $O(\gs^3)$, and the quintic term starts at $\gs^3$.

\emph{(ii) The $\Sigma_2$-equation at quartic order.}
The truncated equation reads $\vertex{4}(\sigma_2, \Psi^{* 2} - \Sigma_2^*) = 0$ for all $\sigma_2 \in \mathcal U^{\wedge 2}$, and since $\Psi^{* 2} - \Sigma_2^* \in \mathcal U^{\wedge 2}$ the non-degeneracy~\eqref{p:eq:trunc:nondegenerate} gives:
\begin{equation}
\label{p:eq:trunc:saddle-quartic}
\Sigma_2^*
  = \Psi^{* 2},
\end{equation}
as in the untruncated theory and without any use of the universality assumption.
The full equation is $\vertex{4}(\sigma_2, \Psi^{* 2} - \Sigma_2^*) = 0$ for all $\sigma_2 \in \Hn{2}$, and it holds at~\eqref{p:eq:trunc:saddle-quartic} for every $\sigma_2$ since its argument vanishes.
Together with (i), the truncated solution satisfies every component of the full equations of motion of $\caction[\le 4]$ except those of $\Psi$ along the universal directions outside $\mathcal U$.

\emph{(iii) The $\Sigma_2$-equation at quintic and sextic order.}
Along $\sigma_2 \in \mathcal U \wedge \mathcal N$, the terms $\vertex{4}(\sigma_2, \Psi^2 - \Sigma_2)$ and $\vertex{5}(\sigma_2, \Sigma_2, \Psi)$ of~\eqref{p:eq:trunc:eom-sigma} are, after expanding $\Sigma_2$, sums of vertices with exactly one non-universal leg.
They vanish identically by the universality assumption.
At quintic order this is the whole equation, which therefore holds identically along $\mathcal U \wedge \mathcal N$.
Along the other pairs outside $\mathcal U^{\wedge 2}$ the assumption gives nothing.
To find the residual at leading order we solve the truncated equation perturbatively.
We write $\Sigma_2^{(n)*[\mathcal U]}$ for the solution of the truncated equation of $\Sigma_2$ of $\caction[\le n+2]^{[\mathcal U]}$ at fixed $\Psi \in \mathcal U$, on the branch connected to $\Psi^2$, and $\Sigma_2^{(n)*}$ for the solution of the full one.
The component $\Sigma_2^*$ of a solution of the truncated equations is $\Sigma_2^{(n)*[\mathcal U]}$ evaluated at $\Psi = \Psi^*$.
At quintic order, we write the truncated solution as:
\begin{equation}
  \label{p:eq:trunc:expansion}
  \Sigma_2^{(3)*[\mathcal U]}
    = \Psi^2 + \frac{\gs}{5} \, Y_2^{[\mathcal U]} + O(\gs^2),
  \qquad
  Y_2^{[\mathcal U]} \in \mathcal U^{\wedge 2},
\end{equation}
and the order $\gs$ of the truncated equation along $B \in \mathcal U^{\wedge 2}$ is:
\begin{equation}
  \label{p:eq:trunc:saddle-quintic}
  \vertex{4}(B, Y_2^{[\mathcal U]})
    = \vertex{5}(B, \Psi^2, \Psi)
  \quad
  \forall B \in \mathcal U^{\wedge 2},
\end{equation}
which determines $Y_2^{[\mathcal U]}$ uniquely by~\eqref{p:eq:trunc:nondegenerate}.
In the untruncated theory the same relation with $B$ running over all of $\Hn{2}$ gives the first correction $\sprod{2}{2}^{-1} \circ \sprod{3}{2}(\Psi^3)$.
The truncated $Y_2^{[\mathcal U]}$ is its analogue with the inverse taken on $\mathcal U^{\wedge 2}$.
We insert the expansion into~\eqref{p:eq:trunc:eom-sigma} along any $\sigma_2 \in \Hn{2}$.
The order $\gs^0$ vanishes because $\Sigma_2 - \Psi^2 = O(\gs)$, and the order $\gs$ is:
\begin{equation}
\label{p:eq:trunc:residual}
\pair[2]{\sigma_2}{\mathcal E^{(3)}_{\Sigma_2}}
  = \frac{\gs}{5} \Big[
      \vertex{5}(\sigma_2, \Psi^2, \Psi)
      - \vertex{4}(\sigma_2, Y_2^{[\mathcal U]})
    \Big]
    + O(\gs^2),
\qquad
\sigma_2 \in \Hn{2}.
\end{equation}
The bracket vanishes on $\mathcal U^{\wedge 2}$ by~\eqref{p:eq:trunc:saddle-quintic}.
Outside $\mathcal U^{\wedge 2}$, it is that relation evaluated outside the space on which it was imposed, and it need not vanish.
It vanishes for every $\sigma_2$ exactly when the untruncated correction $\sprod{2}{2}^{-1} \circ \sprod{3}{2}(\Psi^3)$ lies in $\mathcal U^{\wedge 2}$.
Along $\mathcal U \wedge \mathcal N$ it vanishes identically, each of its two terms having a single non-universal leg.
Without the universality assumption, it need not vanish there.

At sextic order the component of $\mathcal E^{(4)}_{\Sigma_2}$ along the same pairs has the same order-$\gs$ term as~\eqref{p:eq:trunc:residual}, because the sextic term moves the solution only at order $\gs^2$.
Along $\mathcal U \wedge \mathcal N$, however, the term $\dvertex{6}(\sigma_2, \Sigma_2^2)$ of~\eqref{p:eq:trunc:eom-sigma} need not vanish.
Its elementary part $\vertex{6}(\sigma_2, \Sigma_2^2)$ does, having a single non-universal leg.
Its exchange part need not: the inverse of $\sprod{2}{2}$ on $\Hn{2}$ connects $\mathcal U \wedge \mathcal N$ to the pairs of two non-universal states through components of $\vertex{4}$ with three non-universal legs, which the assumption allows.
With $\Sigma_2 = \Psi^2 + O(\gs)$ the residual is:
\begin{equation}
\label{p:eq:trunc:residual-sextic}
\begin{aligned}
  \pair[2]{\sigma_2}{\mathcal E^{(4)}_{\Sigma_2}}
    &
    = \frac{\gs^2}{20} \, \dvertex{6}(\sigma_2, \Psi^2, \Psi^2)
      + O(\gs^3)
    \\
    &
    = - \frac{3 \gs^2}{50} \, \vertex{5 \cdot 5}(\sigma_2, \Psi^2, \Psi^2)
      + O(\gs^3),
  \qquad
  \sigma_2 \in \mathcal U \wedge \mathcal N,
\end{aligned}
\end{equation}
where only the exchange part survives in the second equality.
The lemma makes no claim about the other pairs outside $\mathcal U^{\wedge 2}$ beyond order $\gs$, nor about a branch of solutions not connected to $\Psi^2$.

\emph{(iv) The elimination.}
At quartic order, $\Psi^2 \in \mathcal U^{\wedge 2}$ and we can complete the square of \cref{p:sec:hs:quartic} inside the truncation:
\begin{equation}
  \label{p:eq:trunc:eliminate-quartic}
  \caction[\le 4]^{[\mathcal U]}
    = S_{\le 4} \big|_{\mathcal U}
      - \frac{\gs^2}{4!} \, \vertex{4}(\Sigma_2 - \Psi^2, \Sigma_2 - \Psi^2),
  \qquad
  \caction[\le 4]^{[\mathcal U]} \Big|_{\Sigma_2 = \Psi^2}
    = S_{\le 4} \big|_{\mathcal U},
\end{equation}
with no remainder: the truncated elimination is exact, as the full one is.
This is the step that fails for an auxiliary space not containing $\Psi^2$, where the square cannot be completed and the eliminated action differs from $S_{\le 4}$ at order $\gs^2$ for generic vertices.

At quintic order the truncated action is quadratic in $\Sigma_2$, and we evaluate it on the expansion~\eqref{p:eq:trunc:expansion} of the truncated solution $\Sigma_2^{(3)*[\mathcal U]}$.
Through order $\gs^4$ only $Y_2^{[\mathcal U]}$ is needed.
The terms linear in $\Sigma_2 - \Psi^2$ from the two quartic terms of~\eqref{p:eq:trunc:action-explicit} cancel each other, and the linear term from the quintic vertex, $\tfrac{\gs^3}{60} \vertex{5}(\Sigma_2 - \Psi^2, \Psi^2, \Psi)$, feels the order-$\gs^2$ term of the solution only at order $\gs^5$.
The orders $\gs^2$ and $\gs^3$ reproduce $\vertex{4}(\Psi^4)$ and $\vertex{5}(\Psi^5)$ with the coefficients of $S_{\le 5}$.
At order $\gs^4$ the linear term of the quintic vertex and the completed square of~\eqref{p:eq:trunc:eliminate-quartic} give:
\begin{equation}
  \label{p:eq:trunc:eliminate-quintic}
  \begin{aligned}
    \caction[\le 5]^{[\mathcal U]} \Big|_{\Sigma_2^{(3)*[\mathcal U]}}
      &
      = S_{\le 5} \big|_{\mathcal U}
        + \gs^4 \Big[
          \frac{1}{300} \, \vertex{5}(Y_2^{[\mathcal U]}, \Psi^2, \Psi)
          - \frac{1}{600} \, \vertex{4}(Y_2^{[\mathcal U]}, Y_2^{[\mathcal U]})
        \Big]
        + O(\gs^5)
      \\
      &
      = S_{\le 5} \big|_{\mathcal U}
        + \frac{\gs^4}{600} \, \vertex{5}(Y_2^{[\mathcal U]}, \Psi^2, \Psi)
        + O(\gs^5),
  \end{aligned}
\end{equation}
where the second line uses~\eqref{p:eq:trunc:saddle-quintic} with $B = Y_2^{[\mathcal U]}$.
The value is a rational function of $\Psi$, with denominator the determinant of the dressed kinetic operator $\kin{2}{\Psi}{3}$ restricted to $\mathcal U^{\wedge 2}$, and~\eqref{p:eq:trunc:eliminate-quintic} is its Taylor series.
The elimination in the full theory gives the same expression with $\sprod{2}{2}^{-1} \circ \sprod{3}{2}(\Psi^3)$ in place of $Y_2^{[\mathcal U]}$, so the two agree with $S_{\le 5}|_{\mathcal U}$ through order $\gs^3$, which is (iv) at quintic order.
Their order-$\gs^4$ terms differ in general.

To compare them, we fix a basis of $\Hn{2}$ adapted to the decomposition $\Hn{2} = \mathcal U^{\wedge 2} \oplus (\mathcal U^{\wedge 2})^c$, with $(\mathcal U^{\wedge 2})^c$ the sum of the last two summands of~\eqref{p:eq:trunc:split-H2}.
In this basis we write:
\begin{equation}
  \label{p:eq:trunc:blocks}
  \vertex{4}
    = \begin{pmatrix}
        a & b
        \\
        b^{\mathsf T} & d
      \end{pmatrix},
  \qquad
  \vertex{5}(\,\cdot\,, \Psi^2, \Psi)
    = (v_{\mathcal U}, v_c),
\end{equation}
for the matrix of the quartic form and the components of the source.
The source is the vector of the values $\vertex{5}(B, \Psi^2, \Psi) = \pair[2]{B}{\sprod{3}{2}(\Psi^3)}$ on the basis.
The composite vertex on the diagonal is $\vertex{5 \cdot 5}(\Psi^6) = \vertex{5}(\Psi^2, \Psi, \sprod{2}{2}^{-1} \circ \sprod{3}{2}(\Psi^3))$, and the block inversion of~\eqref{p:eq:trunc:blocks} gives:
\begin{equation}
  \label{p:eq:trunc:schur}
  \vertex{5 \cdot 5}(\Psi^6)
    = v_{\mathcal U}^{\mathsf T} \, a^{-1} v_{\mathcal U}
      + w^{\mathsf T} m^{-1} w,
  \qquad
  m
    := d - b^{\mathsf T} a^{-1} b,
  \qquad
  w
    := v_c - b^{\mathsf T} a^{-1} v_{\mathcal U},
\end{equation}
where $m$ is invertible whenever $a$ and $\sprod{2}{2}$ are.
The quantities $v_{\mathcal U}^{\mathsf T} a^{-1} v_{\mathcal U}$ and $w^{\mathsf T} m^{-1} w$ depend neither on the basis nor on the choice of $\mathcal U^c$: $w$ is the residual~\eqref{p:eq:trunc:residual}, a linear form vanishing on $\mathcal U^{\wedge 2}$, and $m^{-1}$ is the inverse quartic form restricted to such linear forms.
The left side is the untruncated order-$\gs^4$ coefficient of~\eqref{p:eq:trunc:eliminate-quintic}, up to the factor $\gs^4 / 600$.
The first term on the right is the truncated one, since $Y_2^{[\mathcal U]} = a^{-1} v_{\mathcal U}$ by~\eqref{p:eq:trunc:saddle-quintic}.
Hence:
\begin{equation}
  \label{p:eq:trunc:schur-difference}
  \caction[\le 5]^{[\mathcal U]} \Big|_{\Sigma_2^{(3)*[\mathcal U]}}
    - \caction[\le 5] \Big|_{\Sigma_2^{(3)*}}
    = - \frac{\gs^4}{600} \, w^{\mathsf T} m^{-1} w
      + O(\gs^5),
\end{equation}
both at the same $\Psi \in \mathcal U$.
The vector $w$ is, up to the factor $5 / \gs$, the residual~\eqref{p:eq:trunc:residual} along $(\mathcal U^{\wedge 2})^c$, and under the universality assumption it vanishes on $\mathcal U \wedge \mathcal N$.
The difference is thus quadratic in the violated component of the auxiliary equation, and its sign is not fixed, $m$ being indefinite like $\vertex{4}$.
Because the quintic action is quadratic in $\Sigma_2$, this holds exactly, wherever the inverses exist:
\begin{equation}
  \label{p:eq:trunc:exact-difference}
  \caction[\le 5] \Big|_{\Sigma_2^{(3)*}}
    - \caction[\le 5]^{[\mathcal U]} \Big|_{\Sigma_2^{(3)*[\mathcal U]}}
    = \frac{\gs^2}{24} \, e_c^{\mathsf T} \big( (\kin{2}{\Psi}{3})^{-1} \big)_{cc} \, e_c,
\end{equation}
with $e_c$ the component of $\mathcal E^{(3)}_{\Sigma_2}$ along $(\mathcal U^{\wedge 2})^c$ at the truncated solution, whose leading term is~\eqref{p:eq:trunc:residual}.
The block $\big( (\kin{2}{\Psi}{3})^{-1} \big)_{cc}$ is that of the inverse of the dressed kinetic operator on $(\mathcal U^{\wedge 2})^c$.

At sextic order the sextic term contributes $\tfrac{\gs^4}{6!} \dvertex{6}(\Psi^6)$ at order $\gs^4$.
It moves the solution at order $\gs^2$, which changes the value only at order $\gs^6$, since the quintic action is stationary at the quintic solution.
With the reduced vertex of the untruncated theory, $\dvertex{6}(\Psi^6) = \vertex{6}(\Psi^6) - \tfrac{6}{5} \, \vertex{5 \cdot 5}(\Psi^6)$.
Adding~\eqref{p:eq:trunc:eliminate-quintic} and using~\eqref{p:eq:trunc:schur}, we get:
\begin{equation}
  \label{p:eq:trunc:eliminate-sextic}
  \caction[\le 6]^{[\mathcal U]} \Big|_{\Sigma_2^{(4)*[\mathcal U]}}
    = S_{\le 6} \big|_{\mathcal U}
      - \frac{\gs^4}{600} \, w^{\mathsf T} m^{-1} w
      + O(\gs^5).
\end{equation}
The truncation matches one order less than the full theory, and the difference is the same term.
This vertex is moreover not computable from the truncation data, since its exchange part needs the inverse of $\sprod{2}{2}$ on all of $\Hn{2}$.
With the truncation-consistent reduced vertex~\eqref{p:eq:sol:W6-truncated}, whose exchange part $\vertex{5 \cdot 5}^{[\mathcal U]}$ is built with the inverse $(\sprod{2}{2}^{[\mathcal U]})^{-1}$ on $\mathcal U^{\wedge 2}$, the matrix $a^{-1}$ of~\eqref{p:eq:trunc:blocks}, and needs only vertices with all legs in $\mathcal U$, we have:
\begin{equation}
  \label{p:eq:trunc:W6-truncated}
  \vertex{5 \cdot 5}^{[\mathcal U]}(\Psi^6)
    = \vertex{5}(Y_2^{[\mathcal U]}, \Psi^2, \Psi),
  \qquad
  \frac{\gs^4}{720} \, \dvertex{6}^{[\mathcal U]}(\Psi^6)
    + \frac{\gs^4}{600} \, \vertex{5 \cdot 5}^{[\mathcal U]}(\Psi^6)
    = \frac{\gs^4}{720} \, \vertex{6}(\Psi^6),
\end{equation}
the first identity being~\eqref{p:eq:trunc:saddle-quintic} read on $\mathcal U^{\wedge 2}$.
Hence the action~\eqref{p:eq:trunc:action} with $\dvertex{6}$ replaced by $\dvertex{6}^{[\mathcal U]}$ equals $S_{\le 6} \big|_{\mathcal U} + O(\gs^5)$ at the solution of its equation of $\Sigma_2$.
The truncated reduction cancels the exchange term of the truncated theory in the same way as the full reduction cancels the full exchange term, and it restores the match through order $\gs^4$ without enlarging the auxiliary space.
This proves (iv).

The price is that this action is no longer the restriction of the full cubic action.
Its solution violates the equation of $\Sigma_2$ of the full action along $\mathcal U^{\wedge 2}$ itself at relative order $\gs^2$:
\begin{equation}
\label{p:eq:trunc:residual-inner}
\pair[2]{\sigma_2}{\mathcal E^{(4)}_{\Sigma_2}}
  = - \frac{3 \gs^2}{50} \, \big( \vertex{5 \cdot 5} - \vertex{5 \cdot 5}^{[\mathcal U]} \big)(\sigma_2, \Psi^2, \Psi^2)
    + O(\gs^3),
\qquad
\sigma_2 \in \mathcal U^{\wedge 2},
\end{equation}
and this violation and the one of (iii) cancel in the value of the action through order $\gs^4$.
\qed

\subparagraph{Scope of the lemma.}

The lemma is a statement about equations of motion and about the value of the action, not about stability.
The Hessian of the full cubic action at the truncated solution has directions outside $\mathcal U$ and outside $\mathcal U^{\wedge 2}$, and nothing in the universality assumption constrains its sign there.
We discuss it in \cref{p:sec:solutions}.

Without the universality assumption the computations of the proof still apply and give the generic statements.
The equation of $\Psi$ along every direction of $\mathcal U^c$ then has the level-truncation residual~\eqref{p:eq:trunc:level-residual}, in general from order $\gs$.
From quintic order on, the auxiliary equation is in general violated along every pair outside $\mathcal U^{\wedge 2}$, by~\eqref{p:eq:trunc:residual}, the pairs $\mathcal U \wedge \mathcal N$ no longer being protected.
The selection rule and the level of $\mathcal U$ are thus independent sources of error: the first removes the error towards the non-universal states, the second leaves it along the universal states outside $\mathcal U$.

The assumption is weaker than a symmetry.
If the vertices were invariant under a $\mathbb Z_2$ under which the non-universal states are odd and the universal ones even, every coupling with an odd number of non-universal legs would vanish.
The component of (iii) along $\mathcal U \wedge \mathcal N$ would then vanish identically at all three orders.
The universality assumption allows couplings with three non-universal legs, which is why the sextic-order residual~\eqref{p:eq:trunc:residual-sextic} survives under it.

Three limitations remain.
The states of the same levels as $\mathcal U$ that lie outside $\mathcal U$ belong to $\mathcal U^c$ and are universal, so the selection rule says nothing about them.
Since $\mathcal U$ contains all the twist-even universal Siegel-gauge states of its levels, these are the twist-odd ones, which decouple from $\mathcal U$ by the worldsheet twist (checked on the quartic contact terms at zero momentum), and the states outside Siegel gauge.
The kinetic block~\eqref{p:eq:sol:kinetic-block} does not follow from the level argument for them.
  If it fails, the equation of $\Psi$ along them is violated already at order $\gs^0$.

The lemma does not address the gauge structure of the auxiliary field.

Finally, the orders in (iii) and (iv) are Taylor orders at fixed $\Psi$.
In the variables $(\gs \Psi, \gs^2 \Sigma_2)$ the combination $\gs^2 \caction$ does not depend on $\gs$, so at a solution with $\gs \Psi = O(1)$, where a tachyon vacuum sits, the order counting carries no accuracy statement.
The exact statements survive there: the identities (i) and (ii), and the exact form~\eqref{p:eq:trunc:exact-difference} of the quintic-order difference, which is the quantity to report next to the energy of a truncated solution.


